\documentclass[11pt]{article}
\usepackage[T1]{fontenc}
\usepackage{lmodern,microtype}
\usepackage[margin=1in]{geometry}
\usepackage{amsmath,amssymb,amsthm,mathtools}
\usepackage{booktabs,tabularx,array,longtable,enumitem}
\usepackage[dvipsnames]{xcolor}
\usepackage[colorlinks=true,linkcolor=MidnightBlue,citecolor=MidnightBlue,urlcolor=MidnightBlue]{hyperref}
\usepackage{bookmark}
\setlength{\emergencystretch}{2em}
\setcounter{tocdepth}{2}
\setlist{topsep=4pt,itemsep=3pt,parsep=0pt}
\newtheorem{theorem}{Theorem}[section]
\newtheorem{proposition}[theorem]{Proposition}
\newtheorem{lemma}[theorem]{Lemma}
\newtheorem{corollary}[theorem]{Corollary}
\theoremstyle{definition}
\newtheorem{definition}[theorem]{Definition}
\theoremstyle{remark}
\newtheorem{remark}[theorem]{Remark}
\newcommand{\R}{\mathbb R}
\newcommand{\C}{\mathbb C}
\newcommand{\N}{\mathbb N}
\newcommand{\dd}{\,\mathrm d}
\newcommand{\ee}{\mathrm e}
\newcommand{\ii}{\mathrm i}
\newcommand{\abs}[1]{\left|#1\right|}
\newcommand{\norm}[1]{\left\lVert#1\right\rVert}
\newcommand{\ip}[2]{\left\langle#1,#2\right\rangle}
\newcommand{\sinc}{\operatorname{sinc}}
\DeclareMathOperator{\Log}{Log}
\DeclareMathOperator{\meas}{meas}
\numberwithin{equation}{section}
\newcolumntype{Y}{>{\raggedright\arraybackslash}X}
\renewcommand{\arraystretch}{1.14}
\title{Boundary Bootstrap in the Newman Flow\\
\large Exterior Zero Enclosures, Local Green Moments,\\
and the Limits of Sparse Phase Alignment}
\author{Edward Baker\thanks{Drafted for Edward Baker with substantial OpenAI Codex LLM assistance.
The analytical input and review provenance are stated in Section~\ref{sec:evidence}.
This is a working research manuscript, not a claim of independent specialist validation.}}
\date{13 November 2026}
\hypersetup{pdftitle={Boundary Bootstrap in the Newman Flow},pdfauthor={Edward Baker},
pdfsubject={Analytical synthesis through GNP-BB-E11}}
\begin{document}
\maketitle
\begin{abstract}
We develop a boundary-bootstrap approach to the canonical Newman heat flow.
A complete corrected finite-sum expansion, an independently nonzero zeta
carrier, and five fixed classical exponent pairs give the following
exterior enclosure: for some unevaluated $t_*>0$, every zero $x+\ii y$ of
$H_t$ with $0<t\le t_*$ and $|x|\ge4\pi\ee^{1/t}$ satisfies
$|y|<567/1000$. The proof first treats a finite window in
$t\log(x/(4\pi))$, then joins a published all-real far-tail theorem.
A separate Green-function argument reduces the deeper width $9/16$,
on the curve $t=1/\log(X/(4\pi))$, to an explicit still-unproved
Poisson moment bound at every center. We localize that moment to a
middle band, prove quantitative center means, and show that a specified
class of growing positive Dirichlet-polynomial resonators has vanishing
normalized whole-dyadic mean. On the shorter triangular window of length
scale $U^{1/32}$, we retain every pair in an exact frozen quadratic form.
Sparse logarithmic packets have a vanishing effect on the complete
moment; an explicit quartic alignment at actual heights has energy
of order $U^{-241/640}$ and complete effect $O(U^{-59/320})$.
These results delimit specific methods without deciding the outstanding
every-center bound or its proposed short-window counterexample.
All deductions state their analytic inputs; finite arithmetic checks
serve only as supporting evidence. The final section identifies the
remaining arithmetic, carrier, compact-height and time-propagation tasks.
\end{abstract}
\clearpage
{\small\tableofcontents}
\clearpage
\section{The question and the analytical structure}\label{sec:framework}

The purpose of the boundary bootstrap is to improve a known enclosure
of the zeros by observing the field just outside a proposed smaller
enclosure. Two kinds of information can support such a step. A uniform
relative approximation to a nonvanishing carrier gives direct zero
exclusion. A smaller amount of boundary information can suffice through
Green's formula, provided that the value at the observation point has an
independent lower bound. The second route keeps cancellations that the
first route may discard.

This manuscript develops the results of the height and time dependent
boundary investigation through its E11 review. Its starting analytic
input is the corrected complex Riemann--Siegel estimate established in
the repository's heat-density investigation. We state the input precisely
in Section~\ref{sec:expansion}, retain its review status, and derive the
positive-time extension needed here. The other imported inputs are the
classical logarithmic-phase exponent-pair estimate and two published
heat-flow zero theorems. The resulting deductions concern the complete
canonical field, not a truncated theta model.

\subsection{The field and the three time conventions}

For real $t$ and complex $z$, set
\begin{equation}\label{eq:canonical}
 H_t(z)=\int_0^\infty \ee^{tu^2}\Phi(u)\cos(zu)\dd u,
 \qquad
 \Phi(u)=\ee^u\sum_{n\ge1}
 (2\pi^2n^4\ee^{8u}-3\pi n^2\ee^{4u})\ee^{-\pi n^2\ee^{4u}}.
\end{equation}
The normalization, derived below, is
\begin{equation}\label{eq:normalization}
 H_0(z)=\frac18\xi\left(\frac{1-\ii z}{2}\right),\qquad
 \partial_tH_t=-\partial_z^2H_t.
\end{equation}
Thus $|x|/2$ is the ordinary zeta ordinate when $z=x+\ii y$.
Throughout, $\log$ is real on positive arguments and $\Log$ is the
branch on the right half-plane continued from the positive axis.
We use
\[
 L_X=\log\frac{X}{4\pi},\qquad t_X=L_X^{-1},\qquad
 \tau_X=\frac X2+\frac{\pi}{8L_X}.
\]
Three different assertions must be distinguished. In a spatial contour
at center $X$, the field is $H_{t_X}$ at one fixed time. In a center
average, $X$ varies and the family follows the curve $t_X$. In a
common-time exterior theorem, one fixed $t$ is used at every height.
The function $z\mapsto H_{t_{\Re z}}(z)$ is not used as a holomorphic
field in a contour identity.

\subsection{Main conclusions and the unresolved premise}

The strongest acquired enclosure is the following quantitative tail
statement. The word ``acquired'' means proved from the named inputs in
this manuscript and its source notes, with the internal reviews described
in Section~\ref{sec:evidence}.

\begin{theorem}[Exterior enclosure]\label{thm:main-width}
There is a finite, unevaluated $t_*>0$ such that
\[
 0<t\le t_*,\quad |x|\ge4\pi\ee^{1/t},\quad H_t(x+\ii y)=0
 \quad\Longrightarrow\quad |y|<\frac{567}{1000}.
\]
There is also a finite onset $X_*$ such that
$H_{t_X}(X+\ii y)=0$, $X\ge X_*$, implies the same strict width.
\end{theorem}

The proof in Section~\ref{sec:width} yields a relative error at most
$1/64$ on its full comparison rectangles. The width improves the
earlier $4/7$ by $31/7000$, at its own onset. It does not specify the
true smallest onset, and it gives no exclusion in the omitted compact
region.

For the Green route, define
\begin{equation}\label{eq:raw-tail-intro}
 \begin{split}
 T_X(v)&=\sum_{\lfloor\ee^{L_X/32}\rfloor<n\le
                    \lfloor\ee^{L_X/2}\rfloor}
 w_n(X)\ee^{\ii(\tau_X+v/2)\log n},\\
 w_n(X)&=n^{-65/64}\exp\left(\frac{(\log n)^2}{4L_X}\right),\qquad
 Q(X)=\int_\R P(v)|T_X(v)|^2\dd v,
 \end{split}
\end{equation}
where the positive strip kernel $P$ is specified in
Section~\ref{sec:moment} and has mass $2/5$.

\begin{proposition}[Conditional deeper step]\label{prop:intro-green}
If $Q(X)\le4/66^2$ for every sufficiently large real $X$, then every
zero $X+\ii y$ of $H_{t_X}$ at sufficiently large $X$ satisfies
$|y|<9/16$.
\end{proposition}

The proposition is a sufficient arithmetic reduction. Its premise
remains unproved, and no converse is asserted. The unconditional
moment estimates obtained below are
\[
 \|T_X-M_X\|_P\ll X^{-1/8192},\qquad
 \frac1U\int_U^{2U}Q(X)\dd X\ll U^{-25/1024},
\]
where $M_X$ is the middle band $3/64<\log n/L_X\le5/16$.
For every interval $I\subset[U,2U]$ of length $h\ge1$,
\[
 \frac1h\int_I Q(X)\dd X
 \ll U^{-1/4096}+U^{-21/512}+\frac{U^{5/128}\log U}{h}.
\]
These estimates become small on specified growing windows. They do not
bound the moment at every center. Section~\ref{sec:means} proves the
stronger, but still averaged, statement for positive resonator weights
of integer support $R\le U^\delta$, for each fixed $0<\delta<123/128$.

At the shorter scale $h_U=U^{1/32}$, define the probability measure
\[
 \dd\nu_{h_U}(s)=h_U^{-1}(1-|s|/h_U)_+\dd s,
 \qquad A_U(Y)=\int Q(Y+s)\dd\nu_{h_U}(s).
\]
An actual witness $A_U(Y)\ge5/66^2$ along an unbounded sequence
would contradict the every-center sufficient budget. It has not been
constructed or excluded. Section~\ref{sec:packets} proves instead that
a specified sparse family of aligned logarithmic packets cannot change
the complete moment by a fixed positive amount. Its explicit quartic
example achieves actual phase alignment but has insufficient mass.

\subsection{Scope and notation}

All implicit constants are uniform in the real centers and in the
fractional parts of the indicated cutoff values. They may depend on
declared fixed strip widths, a fixed finite clock window, a specified
exponent pair, or a fixed parameter written in the subscript.
An onset is chosen after these constants are fixed. An assertion valid
for each fixed width does not imply a uniform onset as that width shrinks.

The classification used here is a quantitative positive-time enclosure,
an explicitly conditional further step, and auxiliary estimates that
exclude named estimating classes. None is counted as a proof of RH or
as an iteration to the real axis. The repository result and avoidance
indices~\cite{RepoResults,RepoAvoid} route respectively to the inherited
claims and the relevant failure patterns. The specific control screen
and the scope of every negative conclusion appear in
Section~\ref{sec:evidence}; this follows the criteria in patterns
01, 02, 07 and 12~\cite{RepoAvoidPatterns}.

\section{The complete field and its finite expansion}\label{sec:expansion}

\subsection{Normalization from the theta primitive}

Let $\vartheta(r)=\sum_{n\in\mathbb Z}\ee^{-\pi n^2r}$ and
$p(u)=\ee^u(\vartheta(\ee^{4u})-1)/2$.
Jacobi reflection~\cite{DLMF} gives $p(-u)-p(u)=\sinh u$, hence $p'(0)=-1/2$.
Direct differentiation shows that $\Phi=(p''-p)/8$, with $\Phi$ as
in~\eqref{eq:canonical}. Applying $\partial_u^2-1$ to the reflection
identity proves that $\Phi$ is even. Its terms are positive on $u\ge0$
because $\pi n^2\ee^{4u}\ge\pi>3/2$.

The complete series and each fixed derivative are bounded by
\[
 C_j\ee^{|u|}\rho(u)^{j+2}\ee^{-\rho(u)},\qquad
 \rho(u)=\pi\ee^{4|u|}.
\]
To verify this, differentiation replaces $2r^2-3r$ by a polynomial
of degree $j+2$, and
$\ee^{-n^2\rho}\le\ee^{-\rho}\ee^{-\pi(n^2-1)}$.
The resulting sum of polynomial coefficients times
$n^{2j+4}\ee^{-\pi(n^2-1)}$ converges. This domination justifies
differentiation of~\eqref{eq:canonical} on compact parameter sets and
proves the heat equation in~\eqref{eq:normalization}.

If $J(z)=\int_0^\infty p(u)\cos(zu)\dd u$, two integrations by
parts give $8H_0(z)=1/2-(1+z^2)J(z)$. Splitting the theta Mellin
integral at one yields
\[
 \xi(s)=\frac12+\frac{s(s-1)}2\int_1^\infty
 \frac{\vartheta(r)-1}{2}
 \{r^{s/2-1}+r^{(1-s)/2-1}\}\dd r.
\]
Substitute $r=\ee^{4u}$ and $s=(1-\ii z)/2$. The integral is
$8J(z)$, and $4s(s-1)=-(1+z^2)$, proving the normalization.
The reflected argument gives the same field. In particular $H_t$ is
even and is real on the real axis for real $t$.

\subsection{The inherited zero-time estimate}

The analysis below uses the following input, with its original status:
``proved in-note from named classical inputs; same-model internal
assessment only''~\cite{HDAComplex,HDAAssessment}. The initial complex
estimate is not re-certified by its use here. It is stronger than
quoting a real-axis Riemann--Siegel formula at a complex point.

Put
\[
 M(w)=\frac{w(w-1)}{16}\pi^{-w/2}\Gamma(w/2),\quad
 w_\pm=\frac{1\pm\ii z}{2},\quad
 a(z)=\sqrt{\frac z{4\pi}},\quad
 \mathcal G(z)=[M(w_+)M(w_-)]^{1/2},
\]
with $\mathcal G$ positive on the large real axis. Define
\[
 \mathcal C(p)=\frac{\cos(2\pi(p^2-p)-\pi/8)}{\cos(2\pi p)},\qquad
 \mathcal R_N(z)=(-1)^N\mathcal G(z)a(z)^{-1/2}\mathcal C(a(z)-N).
\]
The apparent poles of $\mathcal C$ are removable: at each simple
denominator zero the numerator also vanishes. On every fixed compact
neighborhood of $[0,1]$, the function and its fixed derivatives are
bounded.

\begin{proposition}[Named initial input]\label{prop:initial-input}
Fix $r,b>0$ and $N=\lfloor\sqrt{x/(4\pi)}\rfloor$.
On $|\Re z-x|\le r$, $|\Im z|\le b$, at sufficiently large $x$,
\begin{equation}\label{eq:initial-RS}
 H_0(z)=\sum_{n\le N}\{M(w_+)n^{-w_+}+M(w_-)n^{-w_-}\}
       +\mathcal R_N(z)+O_{r,b}(x^{-3/4}E_{\Gamma,0}(z)),
\end{equation}
where $E_{\Gamma,0}=|M(w_+)|+|M(w_-)|$.
The cutoff is fixed on the chart, and the bound is uniform in its
fractional position. Sectorial calibration gives the elementary arms
below with relative error $O_{r,b}(x^{-1})$ and logarithmic slopes
$L(z)/2\pm\ii\pi/4+O_{r,b}(x^{-1})$, where
$L(z)=\Log(z/(4\pi))$.
\end{proposition}

For precision, the source of~\eqref{eq:initial-RS} is HDA CS3.
Its proof retains the ordinary Dirichlet tail on the upper edge of
a growing auxiliary rectangle, pays its vertical edges, and uses the
exact recurrence
\[
 \mathcal C(p)+\mathcal C(p-1)
 =2\cos(2\pi(p-1)^2-\pi/8)
\]
to pay cutoff crossings on the real edge. The source is uniform on
every declared fixed enlargement; no fixed-chart constant is applied
to an unannounced growing strip. HDA CS16--17 also supply, on a fixed
base chart, constants $c_b\ge1$ and $P_b=b+12$ such that
\begin{equation}\label{eq:source-ray}
 \frac{|H_0(z-2\ii v)|}{\mathcal G(x)}
 \ll x^{P_b}\exp\{ |v|L_x/2+c_b|v|+c_bv^2/x\}
 \quad (|v|\le x^{1/3}).
\end{equation}
For all $v$, a logarithmic majorant of the same ratio is
\[
 P_b\log(x+|v|)+c_b(|v|+1)+\frac{|v|}{2}\log(x+|v|).
\]
These height-relative ray bounds, the calibration, and the initial
estimate are the inherited interface. The constants precede the
Gaussian windows chosen next.

\subsection{A positive-time finite-window extension}

Define the nonzero elementary factors
\begin{equation}\label{eq:carriers}
 \begin{split}
 \mathcal A(t,z)&=\frac{\pi^{1/4}}{16}z^{7/4}
 \exp\left(-\frac{\pi z}{8}+
       \frac{t}{16}(L(z)^2-\pi^2/4)\right),\\
 \Theta(t,z)&=\frac z4(L(z)-1)+\frac{\pi tL(z)}{16}+\frac{7\pi}{8},
 \qquad C_\pm(t,z)=\tfrac12\mathcal A(t,z)\ee^{\pm\ii\Theta(t,z)}.
 \end{split}
\end{equation}
The shifted zeta argument is
\begin{equation}\label{eq:shifted-s}
 s_t(z)=\frac{1-\ii z}{2}+\frac{tL(z)}4-\frac{\ii\pi t}{8}.
\end{equation}

\begin{proposition}[Complete value expansion]\label{prop:positive-provider}
For each fixed $D>0$ there are finite constants and an onset such that,
for $0\le t\le1/2$, $tL_x\le D$, and $|\Re z-x|\le1$,
$|\Im z|\le9$, with $N=\lfloor\sqrt{x/(4\pi)}\rfloor$ fixed,
\begin{equation}\label{eq:provider}
 \begin{split}
 H_t(z)&=\mathcal A(t,z)\sum_{n\le N}n^{-1/2}
   \exp\left(\frac t4((\log n)^2-L(z)\log n)\right)
   \cos\left(\Theta(t,z)-(z/2+\pi t/8)\log n\right)\\
 &\quad+\ee^{-\pi^2t/64}\mathcal R_N(z)+\mathcal E_N(t,z),\\
 |\mathcal E_N(t,z)|&\le h_Dx^{-3/4}(|C_+(t,z)|+|C_-(t,z)|).
 \end{split}
\end{equation}
The error includes the complete source exterior and all finite modes.
\end{proposition}

\begin{proof}
For $t>0$, absolute integration of the complete theta source and the
Gaussian moment identity give
\begin{equation}\label{eq:gaussian-heat}
 H_t(z)=\int_\R H_0(z-2\ii v)\dd\gamma_t(v),\qquad
 \dd\gamma_t(v)=(\pi t)^{-1/2}\ee^{-v^2/t}\dd v.
\end{equation}
For example, the Gaussian average of $\ee^{2|v|u}$ is at most
$2\ee^{tu^2}$, which is integrable against the complete source.
Fix the base strip $b=9$ and its constants in~\eqref{eq:source-ray},
then put $\mu_*=D/4+1$ and choose
\begin{equation}\label{eq:window-choice}
 V\ge\max\{16,b+8,D,16c_b,2\mu_*,\sqrt{8D(P_b+2)}\}.
\end{equation}
Only after this choice invoke the initial estimate on the fixed
enlargement $|\Re z-x|\le2$, $|\Im z|\le10+2V$.

Write $E_o=|C_+|+|C_-|$ and $b_\pm=L(z)/2\pm\ii\pi/4$.
The exact elementary identity
$C_\pm(t,z)=C_\pm(0,z)\exp(tb_\pm^2/4)$, together with calibration,
implies
\begin{equation}\label{eq:envelope-comparison}
 E_o(t,z)\asymp E_{\Gamma,0}(z)\ee^{tL_x^2/16},
 \qquad E_o(t,z)\gg\mathcal G(x).
\end{equation}
All comparisons are on the reserve target and uniform in $tL_x\le D$.

First pay the true source exterior. For $V<|v|\le x^{1/3}$,
the exponent after inserting~\eqref{eq:source-ray} is
\[
 -v^2/t+|v|L_x/2+c_b|v|+c_bv^2/x\le-v^2/(4t)
\]
once $x\ge4c_b$. The three positive terms are bounded respectively
by $(1/2,1/32,1/8)v^2/t$. Splitting the remaining Gaussian shows that
the normalized integral is
$O(x^{P_b}\ee^{-V^2/(8t)})=O_D(x^{-2})$.
For $q=|v|\ge x^{1/3}$, $x+q\le2q^3$ eventually, and the full-ray
logarithmic majorant is at most $q^2$ for all such $q$, after increasing
the onset. Since $t\le1/2$, the Gaussian leaves a bound
$O(\ee^{-x^{2/3}/(4t)})$. This is smaller than the required power.

On $|v|\le V$, insert~\eqref{eq:initial-RS} at the fixed $N$.
The local initial error integrates to $O(E_ox^{-3/4})$: calibration
bounds its two gamma weights by
$C_{b,V}|M(w_\pm)|\ee^{\pm vL_x/2}$, and both Gaussian moments equal
$\ee^{tL_x^2/16}$.

For an individual exact-gamma mode, put
$m_{\pm,n}=M(w_\pm)n^{-w_\pm}$ and $\alpha=M'/M$. Taylor's theorem
on the fixed enlargement gives
\[
 m_{\pm,n}(z-2\ii v)
 =m_{\pm,n}(z)\ee^{\pm v(\alpha(w_\pm)-\log n)}
      (1+O_{b,V}(v^2/x)).
\]
The quadratic error is independent of $n$, since the mode's
$-w_\pm\log n$ is linear. Set $d_n=\alpha(w_\pm)-\log n$.
Then $|\Im d_n|\ll_b1$ and
$|\Re d_n|\le L_x/2+O_b(x^{-1})$ for $n\le N$.
The tilted real Gaussian has center $\pm t\Re d_n/2$ of magnitude
at most $\mu_*$ and second moment $O_D(1)$; its bounded complex loss
is $\ee^{t(\Im d_n)^2/4}$. Thus the local error is at most
$O_D(x^{-1}|m_{\pm,n}\ee^{td_n^2/4}|)$.
Completing each linear term to the whole line costs additionally
\[
 O\left(|m_{\pm,n}\ee^{td_n^2/4}|
           \ee^{-(V-\mu_*)^2/(2t)}\right).
\]
Calibration replaces the completed mode, to relative $O_D(x^{-1})$,
by
\[
 C_\pm(0,z)n^{-1/2}\ee^{\mp\ii z\log n/2}
          \exp\left(\frac t4(b_\pm-\log n)^2\right).
\]
The exponential change costs $O(t(1+L_x)/x)=O_D(x^{-1})$.

All these errors can be summed. The heat weight after division by
the corresponding $C_\pm(t,z)$ has modulus at most one, since
$\log n\le L_x/2\le\Re L(z)$ at the chosen onset.
For $y\ge0$, the dominant population is bounded by
$\sum_{n\le N}n^{-(1+y)/2}\ll x^{1/4}$.
The other population, with carrier ratio $O_b(x^{-y/2})$, is
$O_b(x^{-y/2}N^{(1+y)/2})\ll_b x^{1/4}$.
Conjugation treats $y<0$. Hence the absolute population is
$O_D(E_ox^{1/4})$; all local Taylor and calibration errors cost
$O_D(E_ox^{-3/4})$. Since $V-\mu_*\ge V/2$,
\eqref{eq:window-choice} pays the entire artificial Gaussian exterior
even after multiplication by $x^{1/4}$.

Finally write $\mathcal R_N=(-1)^N\mathcal G f_N$.
On the fixed enlargement,
\[
 \mathcal G'/\mathcal G=-\pi/8+O(x^{-1}),\qquad
 f_N=O(x^{-1/4}),\qquad f_N'=O(x^{-3/4}),
\]
uniformly in the fractional cutoff. The removable quotient and
$a'(z)=O(x^{-1/2})$ prove the last two bounds. Therefore
\[
 \mathcal R_N(z-2\ii v)=\ee^{\ii\pi v/4}\mathcal R_N(z)
     +O(|\mathcal G(z)|x^{-3/4}|v|\ee^{C_{b,V}|v|})
\]
on $|v|\le V$. The error has bounded Gaussian average, while the
completed linear phase has average $\ee^{-\pi^2t/64}$.
Its artificial exterior is paid by the same window. We never divide
by the correction, which can vanish. This proves~\eqref{eq:provider};
at $t=0$, use~\eqref{eq:initial-RS} directly.
\end{proof}

Every gate in this proof is eventual after the constants are fixed:
$x^{1/3}\ge V$, the initial enlarged-chart onsets, sectorial calibration,
and inequalities such as $C_{b,V}V^2/x\le1/2$.
The source notes~\cite{BB06,BB12} retain their full initial-chart ledger.
The present extension concerns values at nonnegative time only.

\subsection{The exact arithmetic decomposition}

At an evaluation point $z=x+\ii y$, recenter the chart at its real
coordinate $x$ and put $N=\lfloor\ee^{L_x/2}\rfloor$.
Write $L(z)=\ell_z+\ii\beta_z$ and $s_t=\sigma-\ii\tau_D$. Expanding
the cosine in~\eqref{eq:provider} gives
\begin{equation}\label{eq:four-terms}
 \frac{H_t}{C_-}=D_N+R_N+Q_N^{\rm corr}+E_N,
\end{equation}
where
\begin{align*}
 D_N&=\sum_{n\le N} a_D(n)n^{\ii\tau_D},&
 a_D(u)&=u^{-\sigma}\ee^{t(\log u)^2/4},\\
 R_N&=\ee^{2\ii\Theta}\sum_{n\le N}a_R(n)n^{-\ii\tau_R},&
 a_R(u)&=u^{(y-1)/2}\ee^{t((\log u)^2-\ell_z\log u)/4},\\
 \sigma&=(1+y)/2+t\ell_z/4,&
 \tau_D&=x/2+\pi t/8-t\beta_z/4,\\
 \tau_R&=x/2+\pi t/8+t\beta_z/4,&
 Q_N^{\rm corr}&=\ee^{-\pi^2t/64}\mathcal R_N/C_-.
\end{align*}
Here $E_N=\mathcal E_N/C_-$. The small complex-logarithm corrections
are retained in both exact frequencies. Since
$\Theta_z=L(z)/4+\pi t/(16z)$, integration vertically from the real
axis gives, for $y\ge0$,
\begin{equation}\label{eq:ratio-correction}
 |C_+/C_-|\le\ee^{-yL_x/2},\qquad
 Q_N^{\rm corr}=O(x^{-(1+y)/4}),\qquad
 E_N=O_D(x^{-3/4})(1+O(x^{-y/2})).
\end{equation}
For the correction use $\mathcal G\asymp\mathcal A(0,z)/2$,
$a^{-1/2}\mathcal C(a-N)=O(x^{-1/4})$ and
$|\mathcal A(0,z)/\mathcal A(t,z)|\le1$ at large $L_x$.
No floor is differentiated. At a cutoff threshold, the estimates from
the two charts both compare with the same entire field at the same time.

\section{Signed finite sums and the exterior enclosure}\label{sec:width}

This section assembles the finite-pair argument of BB Note18~\cite{BB18}
with the complete positive provider, retaining the restricted A-chain
comparison to explain exactly what the broader pair choice improves.

\subsection{The logarithmic-phase input and its weighted exponent}

The classical exponent-pair estimate used here is the following
specific form~\cite{MontgomeryVaughan}. For a fixed exponent pair
$(k,l)$ and any interval $I\subset[M,2M]$,
\begin{equation}\label{eq:exp-pair}
 \left|\sum_{n\in I}n^{\pm\ii T}\right|
 \ll T^kM^{l-k}+M/T\qquad(T\asymp x).
\end{equation}
The logarithmic model $T\log u/(2\pi)$ has derivative parameter one
in the cited definition; its derivative errors are zero. Conjugation
covers the opposite phase. A singleton costs at most one and is
included in the first term for the pairs used here. Real half-open
endpoints cause no change in the estimate. The secondary $M/T$ term
will always be retained.

Starting with $P=(1/6,2/3)=AB(0,1)$, the operations are
\[
 A(k,l)=\left(\frac{k}{2k+2},\frac{k+l+1}{2k+2}\right),\qquad
 B(k,l)=(l-1/2,k+1/2).
\]
Words act from right to left. Only finitely many fixed pairs enter
each field estimate, so no uniform constant for an infinite sequence
of pairs is assumed.

Put $Y_x=\ee^{L_x}=x/(4\pi)$ and $c=tL_x$.
On a clipped dyadic block with left endpoint $M=Y_x^v$, Abel summation
applies~\eqref{eq:exp-pair} to the positive weights in
\eqref{eq:four-terms}. Indeed
\[
 u\frac{a_D'(u)}{a_D(u)}=-\sigma+\frac t2\log u,\qquad
 u\frac{a_R'(u)}{a_R(u)}=\frac{y-1}{2}-\frac{t\ell_z}{4}
                                      +\frac t2\log u
\]
are bounded on a fixed $c$-window. A bound $J$ for these expressions
gives supremum plus total variation at most
$2^J(1+J\log2)$ times the starting weight.
Since $\ell_z\ge L_x$ at the recentered evaluation point, the dominant
primary power of $Y_x$ is at most
\begin{equation}\label{eq:dominant-power}
 E_{k,l}(v;y,c)
 =k+\left(l-k-\frac{1+y}{2}-\frac c4\right)v+\frac{cv^2}{4}.
\end{equation}
At $c=1$, introduce
\begin{equation}\label{eq:envelope-function}
 g_{k,l}(v)=\frac{k}{v}+l-k-1+\frac v4,
 \qquad E_{k,l}(v;y,1)=v\left(g_{k,l}(v)-\frac{y-1/2}{2}\right).
\end{equation}
This reduces the signed-tail problem to a finite envelope comparison.
Increasing $c$ improves the exponent, because
$\partial_cE=-v(1-v)/4<0$ for $0<v<1$.

\subsection{An exact threshold for the restricted A-chain}

The simpler chain $A^jP$ has
\[
 k_j=(2^{j+3}-2)^{-1},\qquad l_j=1-(j+2)k_j,
 \qquad g_j(v)=k_j/v-(j+3)k_j+v/4.
\]
The formulas follow by induction under $A$.
The affine functions $k_j-(j+3)k_jv$ for adjacent indices $j-1,j$
cross at
\[
 v_j=\frac{2^{j+2}}{(j+1)2^{j+2}+2},\qquad j\ge1.
\]
Their crossings decrease to zero, and the line of index $j$ has the
smaller intercept. Repeated adjacent comparisons therefore place
$g_j$ on the lower envelope over $[v_{j+1},v_j]$, with $g_0$ on
$[v_1,1/2]$. Convexity on each piece reduces its maximum to endpoints.

\begin{proposition}[Scoped A-chain threshold]\label{prop:A-threshold}
For this chain,
\[
 \max_{0<v\le1/2}\min_{j\ge0}g_j(v)=\frac{351}{10304},
 \qquad \gamma:=\frac12+2\frac{351}{10304}=\frac{2927}{5152}.
\]
Every fixed $y_0>\gamma$ up to $7/12$ is available from the associated
finite head and six-pair comparison. The threshold is sharp only for
the stated summable-power head and separate A-chain strict-power
allowances.
\end{proposition}

\begin{proof}
For $0<v\le1/8$, letting $j\to\infty$ gives an infimum at most
$v/4\le1/32$. On $[1/8,1/2]$ use successively the indices
$5,4,3,2,1,0$, with internal joins
$64/385,32/161,16/65,8/25,4/9$.
The consecutive endpoint values are
\[
 \frac1{32},\quad\frac{1663}{49280},\quad\frac{351}{10304},\quad
 \frac{63}{2080},\quad\frac7{400},\quad-\frac1{72},\quad-\frac1{24}.
\]
Each is at most $351/10304$. At $v_*=32/161$, the crossing rule
shows that $g_3(v_*)=g_4(v_*)=351/10304$ is the minimum over all $j$.
This proves the equality.

For $y_0>\gamma$, write $\rho=y_0-\gamma$ and take the absolute head
through $\lfloor Y_x^{1/8}\rfloor$. The weighted head is summable
with a fixed margin, while~\eqref{eq:dominant-power} on the six
pieces is at most $-\rho/16$ at $c=1$. A sufficiently small clock
drift preserves a strict negative power. The complete transfer is the
one proved in the next subsection; its constants and onset may depend
on $\rho$.

For the converse about the allowance, let
$\delta=(y_0-1/2)/2$ and let $a$ be the head cutoff exponent.
A strictly summable power head requires $a<4\delta$, since its
worst exponent is $1+\delta-a/4$. If $1/2<y_0\le\gamma$, then
$4\delta\le4(351/10304)<v_*$, so $v_*$ lies beyond every such head.
The best separate A-chain power there is
$v_*(351/10304-\delta)\ge0$. Moving a block endpoint by
$O(1/L_x)$ changes the envelope by only $O(1/L_x)$ near $v_*$,
because the first derivatives of the displayed powers are uniformly
bounded there. It cannot produce a fixed negative power. At
$y_0\le1/2$ the head and the elementary Euler margin already fail.
This proves the asserted restriction, not a lower bound on the
actual field error. Other pairs and joint cancellation are outside
the restricted class.
\end{proof}

\subsection{Five pairs and complete error accounting}

Set $y_*=567/1000$ and $\kappa=1/200000$.
Use either $1\le c\le C$ at one fixed time, with $C$ fixed, or the
whole rectangle $|x-X|\le1$ at $t=t_X$, with $|c-1|\le10^{-5}$.
In both cases $y_*\le y\le9/8$. The positive provider of
Proposition~\ref{prop:positive-provider} applies at an appropriate
fixed $D>C$. The exact frequencies satisfy $\tau_D,\tau_R\in[x/2,x]$
eventually. Moreover
\begin{equation}\label{eq:carrier-floor-33}
 \sigma\ge\frac{413399}{400000}>1+\frac1{32},\qquad
 |1/\zeta(s_t)|\le\zeta(\sigma)\le1+\frac1{\sigma-1}\le33.
\end{equation}
This follows from the ordinary Euler product in $\Re s>1$.
It is independent of any information about the zeros of $H_t$.

Take $K=\lfloor Y_x^{1/8}\rfloor$ and $N=\lfloor Y_x^{1/2}\rfloor$.
For $n\le K$,
\[
 a_D(n)\le n^{-((1+y)/2+7c/32)}\le n^{-1-1/512}.
\]
Thus $\ee^q-1\le q\ee^q$ removes the heat weight on this finite head
at cost $At$, where
\begin{equation}\label{eq:A-constant}
 A=\frac14\sum_{n\ge2}n^{-1-1/512}(\log n)^2<\infty.
\end{equation}
The ordinary zeta tail beyond $K$ costs at most
$32K^{-1/32}\ll x^{-1/256}$. The heat-weighted series is never
continued to infinity.

For the remaining dominant population use Table~\ref{tab:pairs}.
The exact endpoint values are displayed so the inequality is a finite
rational argument, independent of a floating calculation.
\begin{table}[htbp]
\centering\small
\begin{tabular}{@{}llll@{}}
\toprule
Scale interval $v$ & Pair $(k,l)$ & $g$ at left & $g$ at right\\
\midrule
$[1/8,4/25]$ & $(1/254,247/254)$ & $1/32$ & $841/25400$\\
$[4/25,89/500]$ & $(11/1642,1573/1642)$ & $5443/164200$ & $4883141/146138000$\\
$[89/500,11/50]$ & $(1/102,385/408)$ & $303221/9078000$ & $1873/56100$\\
$[11/50,9/20]$ & $(13/502,445/502)$ & $18371/552200$ & $5531/180720$\\
$[9/20,1/2]$ & $(1/6,2/3)$ & $-37/2160$ & $-1/24$\\
\bottomrule
\end{tabular}
\caption{Five fixed pairs. Their words on $P$ are, in order,
$A^5$, $A^5BA^2$, $A^4BABA$, $A^3BA^2BA$, and the identity.}
\label{tab:pairs}
\end{table}

All entries are at most $q=669/20000$, and $g''(v)=2k/v^3\ge0$.
Since $(y_*-1/2)/2=q+1/20000$ and $v\ge1/8$, the primary power
is at most $-1/160000$ at $c=1$.
The clock drift costs at most $10^{-5}/16=1/1600000$, leaving
$E_{k,l}\le-9/1600000<-\kappa$. For $c\ge1$ there is no loss.
Within each real scale interval use left-open, right-closed dyads,
clipping the last one at the right endpoint. Every integer of
$K<n\le N$ occurs exactly once, and there are $O(1+L_x)$ blocks.
The secondary term costs $O(x^{-1/2})$ on each block because
$a_D(M)\le1$ and $M\ll\sqrt x$.

For the complete reflected sum, including $n=1$, use $P=(1/6,2/3)$
on the clipped dyads from one to $Y_x^{1/2}$. After the carrier
ratio, its primary power is
\[
 E_R(v)=\frac16-\frac y2+\left(\frac y2-\frac c4\right)v
                         +\frac{cv^2}{4},\qquad0\le v\le1/2.
\]
It is convex. Its two endpoints are at most $-701/6000$ and
$-180397/4800000$ in the stated domain, hence below $-1/32$.
Its secondary cost is
$O(x^{-y/2}M^{17/16}/x)=O(x^{-1/2})$, since the reflected
weight is at most $M^{1/16}$.
The correction and the complete remainder have already been paid in
\eqref{eq:ratio-correction}. Adding every contribution proves
\begin{equation}\label{eq:complete-comparison}
 \left|\frac{H_t(z)}{C_-(t,z)}-\zeta(s_t(z))\right|
 \le At+B_D(1+L_x)x^{-\kappa}.
\end{equation}
The constant $B_D$ absorbs finitely many pair constants, variation
factors, block counts and the faster-decaying error terms. It does
not depend on the individual cutoff positions or on $t$ in the
specified window.

\subsection{Onsets, continuum coverage and the far-tail join}

We use two published inputs in the normalization~\eqref{eq:canonical}:
de Bruijn strip contraction as stated in Polymath Theorem~3.2, and
Polymath Theorem~1.5(i)~\cite{Polymath}. From the initial xi strip,
the former gives $|y|\le\sqrt{1-2t}$ for zeros at $0\le t\le1/2$.
The latter supplies an absolute $C_P>0$ such that zeros with
$0<t\le1/2$ and $x\ge\exp(C_P/t)$ are real. We do not set $C_P=1$.

On the curve, $|c-1|\le1/((X-1)L_X)$ throughout $|x-X|\le1$.
The normalized comparison is analytic on a neighborhood of the
closed rectangle, by~\eqref{eq:carrier-floor-33}. Its error is bounded by
\[
 33\{A/L_X+B(1+L_X)(X-1)^{-\kappa}\}.
\]
Beyond a finite onset this is at most $1/64$. Direct nonvanishing on
the rectangle excludes the closed interval $[y_*,9/8]$ at its
centerline. The initial contracted strip and conjugation give the
curve assertion in Theorem~\ref{thm:main-width}. A rectangle is used
at every real center; there is no discretization gap.

For the common-time result, first choose
$C\ge\max(8,C_P)+1$, $D=C+1$, and then all provider and comparison
constants. Let $x_0$ dominate their height gates. Choose $L_*$ so that
\begin{equation}\label{eq:onset}
 \begin{gathered}
 L_*\ge\max\{2,\kappa^{-1},4224A\},\qquad4\pi\ee^{L_*}\ge x_0,\\
 4224B_D(1+L_*)(4\pi\ee^{L_*})^{-\kappa}\le1,
 \qquad t_*=L_*^{-1}.
 \end{gathered}
\end{equation}
The last expression decreases eventually to zero, so this is possible.
If $0<t\le t_*$ and $x\ge4\pi\ee^{1/t}$, then $L_x\ge1/t\ge L_*$.
For $1\le tL_x\le C$, each of the two terms in
\eqref{eq:complete-comparison}, multiplied by 33, is at most $1/128$.
The same direct exclusion applies. For $tL_x\ge C$,
\[
 x\ge4\pi\ee^{C/t}>\ee^{C_P/t},
\]
so the published all-real theorem applies. Both ranges include their
common endpoint. Evenness covers negative $x$. This completes the
proof of Theorem~\ref{thm:main-width}.

The small exponent $\kappa$ and the unevaluated constants make
\eqref{eq:onset} an existence construction, not a practical onset
estimate. The known $4/7$ and $7/12$ results follow at a sufficiently
large common onset; their earlier proofs retain their own constants.
The original one-step gain $1/16$ from the initial boundary is likewise
implied eventually by the stronger fixed width. None of these
implications supplies repeated contraction to zero width.

\section{Geometry, observation and the complete Green defect}\label{sec:geometry}

The preceding proof uses a uniform relative approximation, for which
direct nonvanishing is the strongest immediate consumer. Green's formula
becomes useful when only an integrated boundary estimate is available.
We first separate the geometric facts from the arithmetic information
that must support them. All statements in this section concern one
fixed holomorphic field at one fixed time.

\subsection{A complete signed identity and a quadratic bound}

Let $D$ be a rectangle, $p\in D$, and let $F$ be analytic on a
neighborhood of $\overline D$, with $F(p)\ne0$. Normalize the positive
Green function by a logarithmic pole, so that the Green function of
the unit disk with pole zero is $-\log|z|$. Denote harmonic measure
at $p$ by $\omega_D^p$; it has total mass one.

\begin{lemma}[Harmonic normalization]\label{lem:normalization-green}
If $A$ is analytic and nonzero near $\overline D$, the Green defect
of $H=AF$ equals that of $F$. If there are no boundary zeros, this
defect is exactly
\begin{equation}\label{eq:green-identity}
 \int_{\partial D}\log|F|\dd\omega_D^p-\log|F(p)|
 =\sum_{z_j\in D}m_jG_D(p,z_j).
\end{equation}
\end{lemma}
\begin{proof}
$\log|A|$ is harmonic, so its boundary average equals its value at
$p$. Subtract from $\log|F|$ the finitely many interior logarithmic
poles, with their multiplicities, and solve the remaining harmonic
Dirichlet problem. The logarithmic-pole terms give precisely the
Green functions in~\eqref{eq:green-identity}.
\end{proof}

\begin{lemma}[Quadratic boundary consumer]\label{lem:quadratic-green}
Without excluding boundary zeros, one has
\begin{equation}\label{eq:quadratic-green}
 \sum_{z_j\in D}m_jG_D(p,z_j)
 \le\frac12\log\left(1+
 \frac{\displaystyle\int_{\partial D}|F-1|^2\dd\omega_D^p
                         -|F(p)-1|^2}{|F(p)|^2}\right).
\end{equation}
Consequently, if $|F(p)|\ge m>0$ and a target $K\subset D$ has
$\inf_KG_D(p,z)\ge g>0$, the condition
\begin{equation}\label{eq:quadratic-margin}
 \int_{\partial D}|F-1|^2\dd\omega_D^p<m^2(\ee^{2g}-1)
\end{equation}
excludes every zero in $K$.
\end{lemma}
\begin{proof}
The mean-value property gives $\int F\dd\omega_D^p=F(p)$.
On expanding squares, the linear terms cancel exactly, leaving
\[
 \int|F|^2\dd\omega_D^p
 =|F(p)|^2+\int|F-1|^2\dd\omega_D^p-|F(p)-1|^2.
\]
Map the disk conformally onto the rectangle with its origin at $p$.
On circles of radius less than one, Jensen's formula and concavity of
the logarithm bound the zero sum by half the logarithm of the boundary
mean of $|F|^2$, minus $\log|F(p)|$. As the radius increases to one,
the zero sums increase to the Green sum, while continuity of $|F|^2$
gives convergence of its boundary means. This proves
\eqref{eq:quadratic-green}, including possible boundary zeros.
A zero in $K$ would contribute at least $g$, contradicting the strict
bound in~\eqref{eq:quadratic-margin}.
\end{proof}

The subtraction of $|F(p)-1|^2$ is exact, not an estimated negative
term. The later application may discard it to obtain a simpler
sufficient inequality. A quadratic boundary bound can be much smaller
than a supremum bound. On the unit disk, for example,
$F_m(z)=1+2((1+z)/2)^m$ has supremum error two and boundary quadratic
error tending to zero by dominated convergence.

\subsection{Explicit contour and observation costs}

Write $k=a-\delta$, $b=a+\eta$, $d=b-k$, and take
\[
 D=\{|x-X|<W,\ \ell<y<u\},\qquad
 p=X+\ii b,\quad K=\{X+\ii y:k\le y\le a\}.
\]
Let $e_-=k-\ell>0$ and $e_+=u-b>0$.
Suppose for this paragraph that $|F-1|\le\varepsilon<1$ on
$\overline D$. Every point between $X+\ii k$ and $p$ has a Cauchy
disk of radius $r=\min(W,e_-,e_+)$. Thus
$|F'|\le\varepsilon/r$ and $|F(p)|\ge1-\varepsilon$.
The segment retains at least half that floor if
\begin{equation}\label{eq:cauchy-cost}
 d\le\frac{(1-\varepsilon)r}{2\varepsilon},
 \quad\text{or equivalently}\quad
 \varepsilon\le\frac{r}{2d+r}.
\end{equation}
At fixed vertical height $u-\ell>d$, equal vertical margins maximize
the available radius. This is a geometric optimization under a common
uniform error, not an acquisition of that error.

For a disk of radius $R>d/2$ centered halfway between $p$ and the
deepest target point, its Green function on the diameter gives the
worst target value
\begin{equation}\label{eq:disk-cost}
 g_{\rm disk}=\log\frac{R^2+d^2/4}{Rd}.
\end{equation}
Indeed the signed diameter coordinates of the two endpoints are
$d/2$ and $-d/2$, and the disk formula is
$\log|(R^2-st)/(R(s-t))|$.
It increases as the target moves upward toward the observation.
If $R=d/2+r$ fits inside the rectangle, domain monotonicity gives
$G_D\ge G_{\rm disk}$. The uniform-error logarithmic allowance is
$\log((1+\varepsilon)/(1-\varepsilon))$, so strict exclusion follows if
\begin{equation}\label{eq:disk-error-threshold}
 \varepsilon<\left(\frac{r}{d+r}\right)^2.
\end{equation}
This identity follows by solving the comparison of the two logarithms.
When a uniform error below one is already known on all of $D$, direct
nonvanishing excludes $D$ without either extra condition. Formulas
\eqref{eq:cauchy-cost}--\eqref{eq:disk-error-threshold} quantify the
costs of the alternative consumers and of future weaker inputs.

\subsection{The strip comparison and the cost of the ends}

Set $d_s=u-\ell$, $\alpha=\pi(b-\ell)/d_s$,
$\beta=\pi(y-\ell)/d_s$. Mapping the strip
$S=\{\ell<\Im z<u\}$ to the upper half-plane gives
\begin{equation}\label{eq:strip-green}
 G_S(p,x+\ii y)=\frac12\log
 \frac{\cosh(\pi(x-X)/d_s)-\cos(\alpha+\beta)}
      {\cosh(\pi(x-X)/d_s)-\cos(\alpha-\beta)}.
\end{equation}
On its centerline the minimum over $K$ is at $y=k$, with value
\[
 \log\frac{\sin(\pi(b+k-2\ell)/(2d_s))}
                {\sin(\pi(b-k)/(2d_s))}.
\]
For the rectangle with the same horizontal edges,
\begin{equation}\label{eq:strip-rectangle}
 0\le G_S(p,z)-G_D(p,z)
 \le\{cosh(\pi W/d_s)-1\}^{-1}.
\end{equation}
After cancelling the common pole, the difference is harmonic, vanishes
on the horizontal edges and is bounded on the vertical edges by the
right side, using $\log(1+q)\le q$. The maximum principle proves
the comparison. Enlarging a domain thus improves its Green signal,
but also changes the boundary data that must be paid.

The positive harmonic barrier
\[
 \frac{\cosh(\kappa_s(x-X))\cos(\kappa_s(y-y_m))}
 {\cosh(\kappa_s W)\cos(\kappa_s d_s/2)},
 \qquad y_m=(u+\ell)/2,\quad0<\kappa_s<\pi/d_s,
\]
is at least one on the two vertical sides and nonnegative on the
horizontal edges. Therefore their total harmonic measure at $p$ is
at most $\operatorname{sech}(\kappa_s W)/\cos(\kappa_s d_s/2)$.
The same barrier bounds the ends of a graph domain contained in this
strip. A passage to an infinite graph contour additionally requires
integrable boundary logarithms and a sequence of admissible end cuts.
For example, polynomial bounds on the absolute logarithms are absorbed
by this exponential end estimate. Geometry alone does not supply those
canonical bounds.

\subsection{Observation placement and a conditional scale rule}

With $D$ and $K$ fixed and the observation above $K$ on the centerline,
raising the observation decreases the worst Green signal. To see this,
map the symmetric domain to the disk so that its centerline maps
monotonically to a diameter. If $t(b)>t(y)$ are the real disk
coordinates, the pseudohyperbolic distance
$(t(b)-t(y))/(1-t(b)t(y))$ increases with $b$. Its negative logarithm
is the Green function. The Cauchy radius also decreases, while the
distance to the target increases. Under a common uniform error, the
smallest independently admissible observation buffer is therefore best.
With unequal boundary costs the choice can differ: the lower and upper
strip masses are $(u-b)/d_s$ and $(b-\ell)/d_s$.

A useful scale check is
\[
 k=3a/4,\quad b=9a/8,\quad\ell=a/2,\quad u=11a/8,
 \quad W\ge7a/16.
\]
Then $d=3a/8$, $r=a/4$, $R=7a/16$.
The Cauchy condition is $\varepsilon\le1/4$; the disk signal is
$\log(29/21)$ and its strict error threshold is $\varepsilon<4/25$.
The length and derivative costs cancel under rescaling. Conversely,
if a vertical margin is $o(d)$, these particular thresholds degenerate.
Even the exact Green signal is at most
\[
 \min\{\log(1+2e_-/d),\ \log(1+2e_+/d)\},
\]
by comparison with the two containing half-planes at the deepest point.

A canonical estimate with the same margin on every one of these
scaled collars would support $a_{n+1}=3a_n/4$. No such estimate is
proved here. If its onset were $X_n\to\infty$, only finitely many
steps would be justified at a fixed height. Both the carrier and the
arithmetic error must therefore be controlled uniformly enough for an
actual iteration.

\section{Reduction of the deeper step to a local Poisson moment}\label{sec:moment}

We now derive the conditional reduction of BB Note15~\cite{BB15},
with every part of the boundary included.

\subsection{A collar with a fixed Green margin}

Fix $X$ large and hold $t=t_X$ throughout this section. Set
\[
 \ell=17/32,\quad u=39/64,\quad b=37/64,\quad a=73/128,
 \quad d_s=u-\ell=5/64,\quad W=\frac{3d_s}{\pi}\log X,
\]
\[
 D_X=\{|\Re z-X|<W,\ \ell<\Im z<u\},\quad
 p=X+\ii b,\quad K_X=\{X+\ii y:9/16\le y\le a\}.
\]
The already acquired centerline enclosure is below $a$.
The disk centered at $X+\ii a$ of radius $5/128$ lies in $D_X$
eventually. In units $1/128$, the observation and deepest target have
coordinates $+1,-1$ in a disk of radius five. Formula
\eqref{eq:disk-cost} and domain monotonicity give
\begin{equation}\label{eq:green-signal}
 \inf_{z\in K_X}G_{D_X}(p,z)\ge\log(13/5).
\end{equation}

Normalize by the same independently nonzero carrier,
\[
 F(z)=\frac{H_{t_X}(z)}{C_-(t_X,z)\zeta(s_{t_X}(z))}.
\]
At every point of the closed rectangle recenter the fixed provider
chart at its real coordinate. The time remains $t_X$.
Since $|t_XL_x-1|=O(W/(XL_X))$, an eventual gate gives
\begin{equation}\label{eq:carrier-floor-66}
 \Re s_{t_X}\ge65/64-1/40000>1+1/65,
 \qquad |1/\zeta(s_{t_X})|\le66.
\end{equation}
The denominator is analytic and nonzero on a neighborhood of the
rectangle. At the observation and on the top, the complete comparison
of Section~\ref{sec:width} applies with the same small clock drift:
\begin{equation}\label{eq:anchor}
 |F(p)|\ge15/16,\qquad F=1+o(1)\quad\hbox{uniformly on the top}.
\end{equation}
Only the top and observation use that estimate; it is not assumed
on the deeper bottom edge.

The quadratic consumer excludes $K_X$ if its complete boundary
payment is at most $9/2$, because
\begin{equation}\label{eq:margin-81}
 (15/16)^2\{(13/5)^2-1\}=81/16>9/2.
\end{equation}
Taking $\kappa_s=\pi/(2d_s)$ in the end barrier shows that both
sides together have mass $O(X^{-3/2})$. The complete provider and
positive finite population give $|H_{t_X}/C_-|\ll X^{1/4}$ there.
After~\eqref{eq:carrier-floor-66}, the side quadratic cost is
$O(X^{-1})$. The top costs $o(1)$ by~\eqref{eq:anchor}.
It remains to pay the bottom at its actual harmonic measure.

\subsection{The kernel and the complete bottom-edge reduction}

The lower-edge Poisson kernel of the containing strip at $p$ is
\begin{equation}\label{eq:poisson}
 P(v)=\frac{\sin(3\pi/5)}
 {2d_s\{\cosh(\pi v/d_s)-\cos(3\pi/5)\}},\qquad d_s=5/64.
\end{equation}
It is positive and even. With the Fourier convention
$\widehat P(\omega)=\int P(v)\ee^{\ii\omega v}\dd v$, one has
\begin{equation}\label{eq:poisson-transform}
 \widehat P(\omega)=\frac{\sinh(|\omega|/32)}
                         {\sinh(5|\omega|/64)},\qquad
 \widehat P(0)=\int P(v)\dd v=2/5.
\end{equation}
To derive the transform, extend the lower-boundary datum
$\ee^{\ii\omega v}$ harmonically into the strip with zero datum on
the top. Separation of variables gives height factor
$\sinh((u-b)|\omega|)/\sinh(d_s|\omega|)$.
At zero frequency the extension is linear in height. Equivalently,
the substitution $s=\tanh(\pi v/(2d_s))$ evaluates the integral
of~\eqref{eq:poisson} directly. Also $P(v)\ll\ee^{-\pi|v|/d_s}$
for $|v|\ge1$.

For any nonnegative bottom datum on the rectangle, extend it by zero
to the rest of the strip bottom. Its strip extension is nonnegative
on both rectangle sides and dominates the rectangle extension by the
maximum principle. Thus the rectangle's bottom harmonic measure is
dominated by $P(v)\dd v$ on $|v|\le W$.

We now justify the exact tail in~\eqref{eq:raw-tail-intro}.
Split the dominant arm at $K=\lfloor\ee^{L_X/32}\rfloor$.
On the finite head the weighted summability exponent exceeds
$65/64-1/40000-1/128>1+1/256$, so heat removal costs $O(1/L_X)$.
The ordinary zeta tail costs $O(X^{-1/2080})$ using
\eqref{eq:carrier-floor-66}.
At $y=17/32$, the reflected primary power with $(1/6,2/3)$
has endpoints $-19/192$ and $-11/384$ at $c=1$.
The clock drift can be made to cost less than $1/160000$;
the complete reflected sum is consequently
$O((1+L_X)X^{-1/48})$, including its inverse-frequency terms.
The correction and full remainder in~\eqref{eq:ratio-correction}
also tend to zero uniformly.

At $x=X+v$, $|v|\le W$, the natural upper cutoff
$\lfloor\sqrt{x/(4\pi)}\rfloor$ differs from the one at $X$
by at most one eventually. The boundary term costs $O(X^{-57/128})$.
For common terms the ratio to the frozen summand in $T_X(v)$ is
\[
 \exp\left(-\frac{t_X}{4}(L(X+v+\ii\ell)-L_X)\log n\right)
 =1+O(W/X).
\]
The complete absolute population is $O(X^{1/4})$, so the summed
freezing error is $o(1)$. We have proved, uniformly along the bottom,
\begin{equation}\label{eq:bottom-tail}
 F(X+v+\ii\ell)-1
 =\frac{T_X(v)+o(1)}{\zeta(s_{t_X}(X+v+\ii\ell))}.
\end{equation}
All heat-weighted sums are finite, all cutoffs are actual integer
cutoffs, and the entire correction and source remainder have been
included.

If $Q(X)\le4/66^2$, Poisson domination and the triangle inequality
in $L^2(P(v)\dd v)$ give bottom cost at most
\[
 66^2(\sqrt{Q(X)}+o(1))^2\le4+o(1).
\]
Adding the top and both sides gives at most $9/2$ eventually.
Equations~\eqref{eq:green-signal}--\eqref{eq:margin-81} exclude $K_X$.
The acquired centerline enclosure excludes the portion above $a$,
and conjugation covers the lower half-plane. This proves
Proposition~\ref{prop:intro-green}, conditional on its stated
every-center premise.

\subsection{Why signs and exact coefficients remain essential}

Expanding the finite squared norm yields the exact signed form
\begin{equation}\label{eq:Q-pairs}
 Q(X)=\frac25\sum_{K<n\le N}w_n^2
 +2\Re\sum_{K<m<n\le N}w_mw_n\ee^{\ii\tau_X\log(n/m)}
              \widehat P\left(\tfrac12\log(n/m)\right).
\end{equation}
The diagonal is $O(K^{-25/32})=O(X^{-25/1024})$, since
$w_n^2\le n^{-57/32}$. The issue is the signed off-diagonal sum.
For $N/2<m<n\le N$, the weights are comparable to $X^{-57/128}$
and the multiplier is bounded below by a positive constant.
The termwise absolute allowance from those pairs alone is therefore
of order at least $X^{7/64}$. This says nothing about a lower bound
for the signed canonical value.

The same issue appears in a coefficient-independent Gram estimate.
On a dyad $(D,2D]$, the unphased matrix with entries
$\widehat P(\log(n/m)/2)$ has all entries bounded below by a
positive constant. Its all-ones Rayleigh quotient is $\gg D$.
The center phase is a diagonal unitary conjugation, preserving its
operator norm. Multiplying that norm by $\sum w_n^2$ therefore
incurs a scale at least $D^2w_D^2$. If $D=\ee^{aL_X}$, its power
is $a^2/2-a/32$, positive when $a>1/16$.
This is a limitation of the coefficient-independent norm bound.

There is also an exact Gaussian representation of each finite arm.
For $s=65/64-\ii(\tau_X+v/2)$ and
$D_{A,B}(s)=\sum_{A<n\le B}n^{-s}$,
\begin{equation}\label{eq:weight-gaussian}
 \sum_{A<n\le B}w_n(X)\ee^{\ii(\tau_X+v/2)\log n}
 =\sqrt{L_X/\pi}\int_\R\ee^{-L_Xr^2}D_{A,B}(s-r)\dd r.
\end{equation}
This is the real Gaussian moment identity applied to a finite sum.
On a dyad of scale $\ee^{aL_X}$, absolute integration has exponent
$(r-1/64)a-r^2$. Its saddle $r=a/2$ gives $a^2/4-a/64$,
again positive for $a>1/16$. Applying a separate exponent-pair bound
under the integral returns the same weighted block exponent after
square completion. An infinite ordinary series is available only
where $\Re(s-r)>1$. Elsewhere a continuation argument must retain
the exact finite-cutoff remainder; the identity itself supplies no
new cancellation. These scoped comparisons explain why a signed
estimate of the actual moment is still needed.

\section{Localization, center means and positive resonators}\label{sec:means}

The next estimates, from BB Note23~\cite{BB23}, concern the exact $Q$
of~\eqref{eq:raw-tail-intro}.
They do not change its kernel, clock or coefficients. Write
$\|f\|_P^2=\int P(v)|f(v)|^2\dd v$ and define
\[
 I_X^{\rm m}=\{n:\lfloor\ee^{3L_X/64}\rfloor<n\le
                       \lfloor\ee^{5L_X/16}\rfloor\},\qquad
 M_X(v)=\sum_{n\in I_X^{\rm m}}w_n(X)\ee^{\ii(\tau_X+v/2)\log n}.
\]
Put $Q_M(X)=\|M_X\|_P^2$.

\subsection{Uniform removal of the two ends}

\begin{theorem}[Middle-band localization]\label{thm:localization}
Uniformly at every sufficiently large real center,
\begin{equation}\label{eq:localization}
 \|T_X-M_X\|_P\ll X^{-1/8192}.
\end{equation}
In particular $|\sqrt{Q(X)}-\sqrt{Q_M(X)}|$ has the same bound.
\end{theorem}

\begin{proof}
For $K_X<n\le\lfloor\ee^{3L_X/64}\rfloor$,
$w_n(X)\le n^{-1-1/256}$. The absolute sum is
$O(K_X^{-1/256})=O(X^{-1/8192})$, uniformly in $v$.
Its $P$-norm has the same order because $P$ has finite mass.

Cover the high population by clipped dyads starting at
$\ee^{5L_X/16}$ and ending at $\ee^{L_X/2}$, open on the left
and closed on the right. For $|v|\le\log X$, the exact frequency
$\tau_X+v/2$ belongs to $[X/3,2X/3]$ eventually.
Use the fixed pair
\[
 AABABA(1/6,2/3)=(1/24,27/32).
\]
For reference, successive application of $A,B,A,B,A,A$ gives
$(1/14,11/14)$, $(2/7,4/7)$, $(1/9,13/18)$,
$(2/9,11/18)$, $(1/11,3/4)$, $(1/24,27/32)$.
The continuous weight decreases, since
$u\partial_u\log w_u=-65/64+\log u/(2L_X)<0$.
Abel summation on a block $(D,E]$ costs exactly
$w_E+\int_D^E|w_u'|\dd u=w_D$.
Writing $s=\log D/L_X$, its primary exponent is
\[
 E_0(s)=\frac1{24}-\frac{41s}{192}+\frac{s^2}{4},\qquad
 E_0(5/16)=-1/1536,\quad E_0(1/2)=-1/384.
\]
Convexity bounds all of $[5/16,1/2]$ by $-1/1536$.
The secondary term is $O(w_DD/X)=O(X^{-1/2})$.
The sum of the $O(1+L_X)$ blocks is therefore
$O((1+L_X)(X^{-1/1536}+X^{-1/2}))$ on this central $v$-range.

For every real $v$, the full absolute population is at most
$\sum_{n\le N_X}n^{-57/64}\ll X^{7/128}$.
On $|v|>\log X$, exponential decay of $P$ bounds the high population's
squared norm by $O(X^{7/64-\pi/d_s})$.
Taking square roots and adding the three norm payments proves
\eqref{eq:localization}. In particular all Poisson exteriors and
all cutoff crossings are included.
\end{proof}

\subsection{Exact membership intervals and oscillatory integration}

For a unified proof, let $S_X$ denote the raw or middle sum, with
cutoff exponents $(\sigma_0,\rho_0)$ equal respectively to
$(1/32,1/2)$ or $(3/64,5/16)$. Set
\[
 K=\lfloor\ee^{\sigma_0L_U}\rfloor,\quad
 N=\lfloor\ee^{\rho_0L_{2U}}\rfloor,\quad
 \alpha=65/64-\rho_0/4\in(1/2,1).
\]
For an interval $I\subset[U,2U]$ of length $h$, an integer $n$
belongs to $S_X$ precisely on
\begin{equation}\label{eq:membership}
 I_n=I\cap[4\pi n^{1/\rho_0},\,4\pi n^{1/\sigma_0}).
\end{equation}
Thus cutoff endpoints are integrated as endpoints of intervals, not
differentiated. On membership, $w_n(X)\le n^{-\alpha}$ and
\begin{equation}\label{eq:clock-derivatives}
 \frac{w_n'(X)}{w_n(X)}=-\frac{(\log n)^2}{4XL_X^2},\quad
 \tau_X'=\frac12-\frac{\pi}{8XL_X^2}\ge\frac14,\quad
 \tau_X''=\frac{\pi}{8X^2}(L_X^{-2}+2L_X^{-3})>0.
\end{equation}

\begin{lemma}[Pair integration]\label{lem:pair-integration}
For $\lambda\ne0$, on the interval $I_n\cap I_m$,
\begin{equation}\label{eq:pair-integral}
 \left|\int_{I_n\cap I_m}w_n(X)w_m(X)
                   \ee^{\ii\lambda\tau_X}\dd X\right|
 \le\frac{8(nm)^{-\alpha}}{|\lambda|}.
\end{equation}
For $\lambda=0$, the upper bound $h(nm)^{-\alpha}$ applies.
\end{lemma}
\begin{proof}
Let $A=w_nw_m$. By~\eqref{eq:clock-derivatives},
$g=A/\tau_X'$ is positive and nonincreasing. Integration by parts
on a nonempty interval with endpoints $p,q$ gives total cost
\[
 \frac{g(p)+g(q)+\int_p^q|g'(X)|\dd X}{|\lambda|}
 =\frac{2g(p)}{|\lambda|}\le\frac{8(nm)^{-\alpha}}{|\lambda|}.
\]
One-sided values handle membership endpoints. An empty or singleton
interval has zero integral. The zero-frequency statement is immediate.
\end{proof}

\begin{lemma}[Logarithmic harmonic row]\label{lem:harmonic-row}
For nonnegative $b_j$, $1\le j\le J$,
\begin{equation}\label{eq:harmonic-row}
 \sum_{j\ne k}\frac{b_jb_k}{|\log(j/k)|}
 \le3(1+\log J)\sum_j j b_j^2.
\end{equation}
\end{lemma}
\begin{proof}
For a symmetric nonnegative matrix $C$ and weights $p_j>0$,
the inequality $2|u_ju_k|\le|u_j|^2p_k/p_j+|u_k|^2p_j/p_k$ gives
\[
 \sum_{j,k}C_{jk}|u_j u_k|
 \le\sum_j|u_j|^2\sum_kC_{jk}p_k/p_j.
\]
Take $u_j=\sqrt j\,b_j$, $p_j=j^{-1/2}$ and
$C_{jk}=1/(\sqrt{jk}|\log(j/k)|)$ off the diagonal.
The weighted row is $\sum_{k\ne j}1/(k|\log(j/k)|)$.
For $k<j$, its summand is at most $1/k+1/(j-k)$;
for $k>j$ it is at most $1/(k-j)$, by integrating $1/u$ between
the two indices. The three harmonic sums give the stated bound.
\end{proof}

\subsection{Product frequencies and exact coincidences}

Fix a real $R\ge1$ and coefficients with $\sum_{r\le R}|c_r|^2=1$.
They may depend on $U$ but are fixed while $X$ is integrated. Define
\[
 B(X)=\left|\sum_{r\le R}c_r\ee^{\ii\tau_X\log r}\right|^2,
 \qquad D_R=\frac1U\int_U^{2U}B(X)\dd X.
\]
Pair integration with amplitude one, followed by
\eqref{eq:harmonic-row}, gives
\begin{equation}\label{eq:resonator-denominator}
 D_R=1+O(R\log(2R)/U).
\end{equation}

To estimate the weighted moment, group by the integer products
$j=nr$ and put
\[
 b_j=\sum_{\substack{nr=j\\K<n\le N,\ r\le R}}n^{-\alpha}|c_r|,
 \qquad J=N\lfloor R\rfloor.
\]
In the exact squared sum the $X$-phase is
$\tau_X\log(nr/(ms))$; the remaining $v$-factor is
$\ee^{\ii v\log(n/m)/2}$ and has modulus one.
Lemma~\ref{lem:pair-integration} treats every nonzero product gap.
For $nr=ms$, including coincidences with $n\ne m$, use the
zero-frequency bound. Integration against the full Poisson kernel
then gives
\begin{equation}\label{eq:product-master}
 \frac1h\int_I B(X)\|S_X\|_P^2\dd X
 \ll\sum_jb_j^2+\frac{\log(2J)}h\sum_jj b_j^2.
\end{equation}
The two forms on the right require different cutoff information.

Write $r=da$, $s=db$ with $d=\gcd(r,s)$, $\gcd(a,b)=1$.
Then $nr=ms$ means $n=bq$, $m=aq$, and the exact restrictions imply
\[
 K/\min(a,b)<q\le N/\max(a,b).
\]
The matrix entry in the first form is bounded by
\begin{align}
 (ab)^{-\alpha}\sum_{q>K/\min(a,b)}q^{-2\alpha}
 &\ll_\alpha K^{1-2\alpha}(ab)^{-\alpha}\min(a,b)^{2\alpha-1}
 \nonumber\\
 &\le C_\alpha K^{1-2\alpha}\frac{d}{\sqrt{rs}}.
 \label{eq:lower-coincidence}
\end{align}
The power tail estimate also holds when its lower endpoint is below
one, since $2\alpha>1$. For the weighted form $\sum j b_j^2$, use
the upper restriction instead:
\begin{equation}\label{eq:upper-coincidence}
 d(ab)^{1-\alpha}\sum_{q\le N/\max(a,b)}q^{1-2\alpha}
 \ll_\alpha dN^{2-2\alpha}.
\end{equation}
This uses $\alpha<1$; an upper endpoint below one gives an empty sum.

The gcd matrix $G_{rs}=\gcd(r,s)/\sqrt{rs}$ has, with weights
$p_r=r^{-1/2}$, the row estimate
\[
 \sum_{s\le R}G_{rs}\frac{p_s}{p_r}
 =\sum_{s\le R}\frac{\gcd(r,s)}s
 \le\sum_{d\mid r}\sum_{\substack{s\le R\\d\mid s}}\frac ds
 \le d(r)(1+\log R).
\]
For each fixed $\varepsilon>0$, $d(r)\ll_\varepsilon r^\varepsilon$.
One direct proof uses $e+1\le2^e\le p^{\varepsilon e}$ for all
large primes; each of the finitely many remaining primes contributes
a finite maximum of $(e+1)p^{-\varepsilon e}$.
The weighted matrix inequality in Lemma~\ref{lem:harmonic-row},
and $d\le R\,d/\sqrt{rs}$ in~\eqref{eq:upper-coincidence}, now give
\begin{equation}\label{eq:coincidence-forms}
 \begin{split}
 \sum_jb_j^2&\ll_\varepsilon K^{1-2\alpha}R^\varepsilon(1+\log R),\\
 \sum_jj b_j^2&\ll_\varepsilon
          N^{2-2\alpha}R^{1+\varepsilon}(1+\log R).
 \end{split}
\end{equation}
Thus both exact coincidences and nearby distinct products are paid.

\subsection{Center means and their precise ranges}

\begin{theorem}[Mean estimates]\label{thm:means}
For $I\subset[U,2U]$, $1\le h=|I|\le U$,
\begin{align}
 \frac1h\int_I Q_M(X)\dd X
 &\ll U^{-21/512}+\frac{U^{5/128}\log U}{h},\label{eq:middle-mean}\\
 \frac1h\int_I Q(X)\dd X
 &\ll U^{-1/4096}+U^{-21/512}
                        +\frac{U^{5/128}\log U}{h}.\label{eq:raw-short-mean}
\end{align}
The direct raw dyadic estimate is stronger than localization followed
by averaging:
\begin{equation}\label{eq:raw-dyadic}
 \frac1U\int_U^{2U}Q(X)\dd X\ll U^{-25/1024}.
\end{equation}
For each fixed $s_0>0$, the proportion of centers with $Q(X)>s_0$
is $O_{s_0}(U^{-25/1024})$.
\end{theorem}

\begin{proof}
For $R=1$, the two forms in~\eqref{eq:product-master} are bounded
directly by $K^{1-2\alpha}$ and $N^{2-2\alpha}$.
For the middle band, $\alpha=15/16$, $K\asymp U^{3/64}$ and
$N\ll U^{5/16}$; their powers are $-21/512$ and $5/128$.
This proves~\eqref{eq:middle-mean}.
Use~\eqref{eq:localization} in the joint norm for
$h^{-1}\dd X\,P(v)\dd v$ and then square, proving
\eqref{eq:raw-short-mean} with the $U^{-1/4096}$ payment.
For the raw sum directly, $\alpha=57/64$, $K\asymp U^{1/32}$ and
$N\ll U^{1/2}$. At $h=U$ the two terms are
$U^{-25/1024}$ and $U^{-57/64}\log U$; the latter is absorbed.
The exceptional-set estimate follows by integrating $Q\ge s_0$
over that set.
\end{proof}

The short-center mean tends to zero uniformly in interval position
when $h/(U^{5/128}\log U)\to\infty$, for example when
$h\ge U^{5/128+\eta}$ for fixed $0<\eta\le123/128$.
This does not reach bounded windows or the scale $U^{1/32}$.
No zero count is inferred from the exceptional set of this auxiliary
moment.

\subsection{A uniform restriction on whole-dyadic resonators}

\begin{theorem}[Vanishing normalized resonator mean]\label{thm:resonator}
For each fixed $0<\delta<123/128$, uniformly for $R\le U^\delta$
and all coefficient vectors as above,
\begin{equation}\label{eq:resonator-limit}
 \frac{\int_U^{2U}B(X)Q(X)\dd X}
      {\int_U^{2U}B(X)\dd X}=o(1).
\end{equation}
The denominator is positive and comparable to $U$ for all sufficiently
large $U$ in this range.
\end{theorem}

\begin{proof}
Apply~\eqref{eq:product-master} and~\eqref{eq:coincidence-forms}
to the middle band at $h=U$. Then transfer the omitted modes using
the nonnegative joint measure $U^{-1}B(X)\dd X\,P(v)\dd v$.
Their norm is $O(U^{-1/8192}D_R^{1/2})$ by localization.
The triangle inequality followed by squaring gives
\begin{equation}\label{eq:resonator-full}
 \begin{split}
 \frac1U\int_U^{2U}BQ\dd X
 &\ll U^{-1/4096}D_R\\
 &\quad+C_\varepsilon\bigl[
 U^{-21/512}R^\varepsilon(1+\log R)\\
 &\hspace{42mm}+U^{-123/128}R^{1+\varepsilon}
                   (1+\log R)\log(2UR)\bigr].
 \end{split}
\end{equation}
This norm transfer retains all cross terms with the omitted modes.
Set $m_\delta=\min(21/512,123/128-\delta)>0$ and
$\varepsilon=m_\delta/(2\delta)$. Each bracketed power is then
at most $U^{-m_\delta/2}$ before logarithms. By
\eqref{eq:resonator-denominator}, $D_R=1+o(1)$ uniformly, so
the normalized mean is
$O_\delta(U^{-1/4096}+U^{-m_\delta/4})$ after absorbing logarithms.
\end{proof}

For comparison, applying the same argument directly to the raw sum
gives
\[
 \frac1U\int_U^{2U}BQ\dd X
 \ll_\varepsilon U^{-25/1024}R^\varepsilon(1+\log R)
 +U^{-57/64}R^{1+\varepsilon}(1+\log R)\log(2UR),
\]
which gives the smaller sufficient open range $\delta<57/64$.
Neither endpoint is asserted to be optimal.
Theorem~\ref{thm:resonator} excludes a fixed positive lower bound
for this normalized whole-dyadic average along an unbounded sequence.
It leaves exceptional individual centers, shorter or adaptive windows,
$X$-dependent coefficient vectors, other weights, and larger supports
outside its scope.

\section{Short-window freezing and actual phase packets}\label{sec:packets}

This section develops the result of BB Note25~\cite{BB25}, accepted
without mathematical correction in the E11 review~\cite{BBReview11}.

Fix exactly
\begin{equation}\label{eq:short-window}
 \begin{gathered}
 h_U=U^{1/32},\qquad Y\in[U+h_U,2U-h_U],\\
 \dd\nu_{h_U}(s)=h_U^{-1}(1-|s|/h_U)_+\dd s,\qquad
 A_U(Y)=\int Q(Y+s)\dd\nu_{h_U}(s).
 \end{gathered}
\end{equation}
The center measure has mass one. The fixed question is whether
$A_U(Y)\ge5/66^2$ at admissible actual centers along an unbounded
sequence of $U$. Positivity would then force a violation of the
every-center $4/66^2$ budget somewhere in the window.
The results of this section do not decide that question.

\subsection{Freezing in the consumed Hilbert norm}

All norms in this subsection use
$\mathcal H_U=L^2(\dd\nu_{h_U}(s)P(v)\dd v)$.
At the fixed center $Y$, write $I_Y=I_Y^{\rm m}$ and set
\[
 G_Y(s,v)=\sum_{n\in I_Y}w_n(Y)
      \ee^{\ii(\tau_Y+(s+v)/2)\log n},\qquad
 F_U(Y)=\|G_Y\|_{\mathcal H_U}^2.
\]
\begin{theorem}[Uniform square-root transfer]\label{thm:freezing}
For every admissible $Y$, uniformly in the fractional parts of both
real cutoff values,
\begin{equation}\label{eq:freezing}
 \|M_{Y+s}-G_Y\|_{\mathcal H_U}\ll U^{-771/16384},\qquad
 |\sqrt{A_U(Y)}-\sqrt{F_U(Y)}|\ll U^{-1/8192}.
\end{equation}
The frozen form has the exact expansion
\begin{equation}\label{eq:frozen-form}
 F_U(Y)=\sum_{n,m\in I_Y}w_n(Y)w_m(Y)\ee^{\ii\tau_Y\log(n/m)}
 \widehat P\left(\frac{\log(n/m)}2\right)
 \sinc^2\left(\frac{h_U\log(n/m)}4\right),
\end{equation}
where $\sinc q=\sin q/q$ with value one at zero.
\end{theorem}

\begin{proof}
For either cutoff exponent $p\in\{3/64,5/16\}$, the real function
$\ee^{pL_X}$ has derivative $p\ee^{pL_X}/X$. Its excursion across
the whole window is $O(U^{p-31/32})=o(1)$.
Since the function is monotone, each floor can change by at most one
integer. At a changed endpoint, $\log n=pL_Y+O(1)$ and
$L_X-L_Y=O(h_U/U)$, so
\[
 w_n(X),w_n(Y)\ll U^{-65p/64+p^2/4}.
\]
The two endpoint powers are $-771/16384$ and $-75/256$.
This also bounds terms that occur only in the frozen sum.

For common indices,
\[
 \partial_X\log w_n(X)=O(U^{-1}),\qquad
 \tau_{Y+s}-\tau_Y-s/2
 =\frac\pi8(L_{Y+s}^{-1}-L_Y^{-1})
 =O\left(\frac{h_U}{U\log^2U}\right).
\]
After multiplication by $\log n$, the phase error is
$O(h_U/(U\log U))$. On the middle band,
$w_n(Y)\le n^{-15/16}$ and $\sum_{I_Y}w_n(Y)\ll U^{5/256}$.
The total common-term error is therefore
$O((h_U/U)U^{5/256})=O(U^{-243/256})$, uniformly for all real $v$.
Adding the two endpoint terms proves the first norm bound.
Theorem~\ref{thm:localization}, integrated against the probability
measure $\nu_{h_U}$, pays the raw-to-middle norm by $U^{-1/8192}$.
The reverse triangle inequality proves the second bound.

The normalized indicator of $[-h_U/2,h_U/2]$ has Fourier transform
$\sinc(h_U\omega/2)$; its self-convolution has precisely the triangular
density in~\eqref{eq:short-window}. In the finite squared sum,
the $s$ and $v$ frequencies are both $\log(n/m)/2$.
Their transforms give~\eqref{eq:frozen-form}. Every ordered pair is
retained. The form is nonnegative because it is a squared joint norm,
not because its signed cross terms have a fixed sign.
\end{proof}

The transfer is deliberately between square-root moments. It is not
an additive $o(1)$ estimate for the moments themselves. The full norm
bound proved next is $O(U^{1/256})$; multiplying it by the localization
error would give $U^{31/8192}$, which grows. A fixed strict surplus
in the square-root norm can be transferred; an equality-level bound
requires its error payment.

\subsection{Summable logarithmic cells and sparse deletion}

Partition the middle indices into cells
\[
 \mathcal C_j=\{n\in I_Y:j\le h_U\log n<j+1\},\qquad
 a_j=\sum_{n\in\mathcal C_j}w_n(Y).
\]
For $n,m$ in cells $j,k$, the absolute multiplier in
\eqref{eq:frozen-form} is at most $C(1+|j-k|)^{-2}$.
Indeed for $|j-k|\ge2$ one has
$h_U|\log(n/m)|\ge|j-k|-1$ and
$|\sinc q|\le\min(1,|q|^{-1})$; the bound
$|\widehat P|\le2/5$ covers the remaining cases.
Since these rows are summable,
\begin{equation}\label{eq:cell-row}
 F_U(Y)\le C\sum_ja_j^2.
\end{equation}
This follows by summing $2a_ja_k\le a_j^2+a_k^2$.
It applies to any subset of indices and even to arbitrary fixed
coefficient arguments, so it is an upper comparison only.

Put $D_j=\ee^{j/h_U}$. The number of integers in a cell is at most
$D_j(\ee^{1/h_U}-1)+1$. In every occupied cell,
$D_j/h_U\gg U^{1/64}\to\infty$, because the lower middle scale is
$U^{3/64}$. Thus the $+1$ is absorbed and the number is $O(D_j/h_U)$.
Using $w_n\le n^{-15/16}$ gives
\begin{equation}\label{eq:cell-mass}
 a_j\ll D_j^{1/16}/h_U,\qquad
 \max_ja_j^2\ll h_U^{-2}U^{5/128}=U^{-3/128}.
\end{equation}
If $j_*$ is the last occupied cell, the full geometric sum satisfies
\[
 \sum_{j\le j_*}\ee^{j/(8h_U)}
 =\frac{\ee^{j_*/(8h_U)}}{1-\ee^{-1/(8h_U)}}
 \ll h_U U^{5/128}.
\]
Consequently~\eqref{eq:cell-row} and~\eqref{eq:freezing} yield
\begin{equation}\label{eq:all-cell-bound}
 F_U(Y)\ll U^{1/128},\qquad A_U(Y)\ll U^{1/128}.
\end{equation}
The triangular window removes the logarithm from the earlier
short-center allowance, but the new allowance still grows.

\begin{theorem}[Sparse-packet effect on the complete moment]\label{thm:sparse}
Let $S\subset I_Y$ meet at most $r\ge1$ cells.
Its frozen standalone energy is $O(rU^{-3/128})$.
Let $A_U^{\setminus S}(Y)$ denote the actual average after deleting
these fixed indices from the complete raw sum at every center in the
window. Then
\begin{equation}\label{eq:sparse-full}
 |A_U(Y)-A_U^{\setminus S}(Y)|
 \ll\sqrt r\,U^{-1/128}+rU^{-3/128}.
\end{equation}
In particular its complete effect tends to zero for $r=o(U^{1/64})$.
\end{theorem}

\begin{proof}
The standalone frozen bound follows from~\eqref{eq:cell-row} and
\eqref{eq:cell-mass}. The whole fixed set $I_Y$ lies strictly inside
the raw cutoffs throughout the window eventually, because
$1/32<3/64<5/16<1/2$ and $L_{Y+s}/L_Y\to1$ uniformly.
There is no moving membership for the selected packet
\[
 P_{S,X}(v)=\sum_{n\in S}w_n(X)\ee^{\ii(\tau_X+v/2)\log n}.
\]
The common-term freezing proof pays its error by $O(U^{-243/256})$.
Its actual joint norm is therefore $O(\sqrt r\,U^{-3/256})$.
For $V=T_{Y+s}$ and $W=P_{S,Y+s}$ in $\mathcal H_U$,
\[
 \|V\|^2-\|V-W\|^2=2\Re\ip VW-\|W\|^2.
\]
Using the complete bound $\|V\|\ll U^{1/256}$ from
\eqref{eq:all-cell-bound} proves~\eqref{eq:sparse-full}.
This includes the interaction with every low, middle and high index.
\end{proof}

The standalone energy vanishes in the different, larger sufficient
range $r=o(U^{3/128})$. Conversely, a fixed positive standalone energy
would require $r\gg U^{3/128}$ under this estimate.
These statements must not be interchanged with the complete-perturbation
range. None bounds the energy of the untouched remainder below the
witness threshold.

\subsection{Quartic integer alignment at actual heights}

We now test one concrete lower mechanism without replacing the physical
phases by freely chosen arguments. Let $D\to\infty$ through integers,
and choose the unique sufficiently large real $Y_D$ satisfying
\begin{equation}\label{eq:quartic-center}
 \tau_{Y_D}=24\pi D^4.
\end{equation}
Existence and uniqueness on the large branch follow from continuity,
$\tau_X'\ge1/4$ and $\tau_X\to\infty$.
Then $Y_D\sim48\pi D^4$. Set $U=2Y_D/3$, so that the center is
admissible, and keep $h_U=U^{1/32}\asymp D^{1/8}$.
Take the fixed packet
\[
 J=\lfloor D^{1/5}/4\rfloor,\qquad S_D=\{D,D+1,\ldots,D+J\}.
\]
Uniformly across the window, $L_X=4\log D+\log12+o(1)$ and
$\log(D+j)/L_X\to1/4$. Thus the packet belongs to both middle and
raw populations throughout the window, with room to spare at both
integer cutoffs.

\begin{theorem}[An aligned packet with vanishing complete effect]\label{thm:quartic}
For the actual centers~\eqref{eq:quartic-center},
\begin{equation}\label{eq:quartic-energy}
 \|P_{S_D,Y_D+s}\|_{\mathcal H_U}^2
 \asymp D^{-241/160}\asymp U^{-241/640},
\end{equation}
and
\begin{equation}\label{eq:quartic-full}
 |A_U(Y_D)-A_U^{\setminus S_D}(Y_D)|\ll U^{-59/320}.
\end{equation}
\end{theorem}

\begin{proof}
Integrating the geometric-series remainder for $1/(1+x)$ gives
\begin{equation}\label{eq:quartic-taylor}
 \begin{split}
 24\pi D^4\log(1+j/D)
 &=2\pi(12D^3j-6D^2j^2+4Dj^3-3j^4)+R_j,\\
 0\le R_j&\le\frac{24\pi}{5}\frac{j^5}{D}
                  \le\frac{3\pi}{640}<\frac1{32}.
 \end{split}
\end{equation}
The polynomial in parentheses is an integer.
For $|X-Y_D|\le h_U$ and $|v|\le1$, the additional relative phase
against the term $n=D$ is
\[
 (\tau_X-\tau_{Y_D}+v/2)\log(1+j/D)
 =O((h_U+1)J/D)=O(D^{-27/40}).
\]
Eventually all packet terms therefore have relative arguments within
$1/16$ modulo $2\pi$ on this entire portion of the joint domain.
Their weights satisfy $w_{D+j}(X)\asymp D^{-61/64}$, since
$(\log(D+j))^2/(4L_X)=\log D/16+O(1)$ uniformly.
After rotation by the first term's phase, real projection gives a
lower bound $\cos(1/16)\sum_{S_D}w_n(X)$ on $|v|\le1$.
The triangle inequality gives the corresponding upper bound for all
$v$. The restricted Poisson mass is positive and the center measure
has mass one. Hence the energy has order
\[
 (D^{1/5}D^{-61/64})^2=D^{2/5-61/32}=D^{-241/160},
\]
proving~\eqref{eq:quartic-energy}.
Finally apply the exact squared-norm difference used in
Theorem~\ref{thm:sparse}, now with the sharper packet norm:
\[
 |A_U-A_U^{\setminus S_D}|
 \ll U^{1/256}U^{-241/1280}+U^{-241/640}
 \ll U^{-59/320}.
\]
The first factor is the norm of the complete raw sum, so the whole
background interaction is included.
\end{proof}

The packet also meets at most two logarithmic cells eventually,
because $h_U\log(1+J/D)=o(1)$.
It succeeds at actual alignment and fails at fixed mass. This stops
this isolated quartic-packet mechanism and the stated sparse deletion
class. It neither disproves the short-center witness nor establishes
the every-center upper bound. A selected positive squared norm is not
a lower bound for the complete squared norm: the complementary vector
can cancel it, and the complementary energy is still unknown.

\section{Scope, collision controls and supporting evidence}\label{sec:evidence}

\subsection{What each argument establishes}

The exterior theorem is a statement about zeros of the complete field
on a specified positive-time exterior. The Green theorem is a conditional
implication with one explicit unpaid premise. The moment theorems are
estimates of the stated finite quadratic forms. The A-chain threshold,
Gram comparison, whole-dyadic resonator theorem and sparse-packet theorem
each restrict a named method or class. None is a universal obstruction
to Green methods, phase cancellation, or an every-center estimate.

The finite widths $7/12$, $4/7$ and $567/1000$ are successive available
bounds at their stated onsets. The A-chain threshold $2927/5152$ is
already passed by the five-pair construction, which illustrates why a
threshold for one allowance cannot be interpreted as a field barrier.
The moment bound $Q\le4/66^2$ is sufficient for the particular $9/16$
consumer. Even a proved violation of that budget would only defeat
this sufficient budget; it would not refute the zero enclosure or
every alternative signed Green estimate.

\subsection{The canonical premises and the control screen}

The control catalog~\cite{RepoControls} contains sources that share
positivity, reflection properties or finitely many observations with
theta and nevertheless have positive-time collisions. Table~\ref{tab:controls}
identifies which premise of the present canonical interpretation each
class fails. The Green, Hilbert-space and cell inequalities themselves
are general analytic statements and may hold in control models.

\begin{longtable}{@{}>{\raggedright\arraybackslash}p{0.36\textwidth}>{\raggedright\arraybackslash}p{0.59\textwidth}@{}}
\caption{Premises needed beyond shared shape or finite observations.}
\label{tab:controls}\\
\toprule
Control class & Premise not supplied\\
\midrule
\endfirsthead
\toprule Control class & Premise not supplied\\\midrule
\endhead
Euler-factor and translated sources & Unchanged ordinary Dirichlet coefficients and the complete canonical two-arm identity.\\
Weighted-square and finite-tail sources & The full unit coefficient sequence together with full normalized reflection.\\
Finite-prime orbit controls & Complete coefficients and reflection beyond the chosen primes and finite coefficient prefix.\\
Reflected super-Gaussian, smooth shape and positive-primitive sources & The exact normalized integer-square theta arithmetic.\\
Positive heat counter-kernels and exact-threshold models & The completed xi source and both of its arithmetic arms.\\
Atomic, finite-profile and compact continuous controls & The complete canonical source, coefficient sequence and profile beyond matched data.\\
RH-conditional normalized endpoint controls & Ordinary Dirichlet arithmetic; their own conditional premise stays conditional.\\
Fixed-amplitude variable readouts & The actual cosine readout and its Newman evolution.\\
Relaxed positive zero measures & Identification with the actual integer zero divisor.\\
Free or finitely calibrated phase variables & The common logarithmic phase at an actual height, with the stated moving population.\\
\bottomrule
\end{longtable}

Reflection rigidity must retain the conjunction of its hypotheses,
not reflection alone. If a control zero lies in a comparison domain
whose carrier is independently nonzero, its relative deviation from
one is exactly one at the zero. If it lies in a Green target, its
positive multiplicity contributes at least the target Green signal.
Such a zero would contradict a proved strict combined margin on that
box. A disjoint collision-free box may still pass. The screen does not
assert that a particular control event occurs in our asymptotic domains.

\subsection{Proof sources and review provenance}

The source notes and their internal reviews supply the following
crosswalk. The manuscript reorganizes their arguments by dependency;
it does not upgrade their evidence level.

\begin{center}\small
\begin{tabularx}{\textwidth}{@{}>{\raggedright\arraybackslash}p{0.30\textwidth}Y Y@{}}
\toprule
Development & Primary local source & Review status retained\\
\midrule
Complex corrected initial expansion & HDA Ansatz Note04, CS3 and CS16--17 & Same-model internal assessment, HDA Review03.\\
First full comparison and geometry & BB Notes06 and 08 & E01--E03 checked in Review03; geometry is explicitly conditional.\\
Positive finite-window provider & BB Note12 & E06 Review05 accepts the named-input proof.\\
Five-pair exterior enclosure & BB Note18 & E08 Review07 accepts curve and common-time statements.\\
Conditional Green reduction & BB Note15 & Conditional consumer; every-center moment remains unpaid.\\
Localization and resonator means & BB Note23 & E10 Review09 accepts the fixed proof.\\
Short-window and actual packet & BB Note25 & E11 Review11 accepts equations (1)--(11), without mathematical correction.\\
\bottomrule
\end{tabularx}
\end{center}

The BB proof and review workers whose dispatch is recorded were selected
as \texttt{gpt-6.1-sol} with reasoning effort \texttt{ultra}; B04 and
B05 conceptual selections also record \texttt{gpt-6-astra}/\texttt{ultra}
dispatches. The B06 and manuscript lead's actual backend model and
reasoning effort are unavailable in the recorded chat configurations
and are not inferred from those selections. The E11 reviewer was
distinct from the B06 proof author and used shared sources. These are
internal LLM assessments, not independent specialist human reviews.
The manuscript's production checks are self-checks of the synthesis,
separate from the E11 analytical review.

The E11 review~\cite{BBReview11} binds the original Note25 bytes and
checks the norm transfer, both floor changes, summable cell rows,
integer Taylor coefficients, and complete background interaction.
The source proof and issued review request were preserved unchanged.
Its acceptance leaves the $5/66^2$ witness and $4/66^2$ upper budget
open, exactly as in Sections~\ref{sec:moment} and~\ref{sec:packets}.

\subsection{The supporting role of computation}

The conclusions above are analytical. Exact rational records replay
the exponent-pair transformations, endpoint comparisons and powers in
the proofs. Hashes identify the reviewed requests and source files.
They do not prove an analytic inequality. No field values, empirical
onset, fitted asymptotic law, or sampled absence of zeros is used in
any theorem here. The review record for E11 independently matches the
nineteen powers in the earlier packet record and the inherited proof
identities; its analytic assessment remains the written review.

The editable TeX, proof crosswalk and small build record are retained
in the investigation. The generated PDF, render images and compilation
products are stored outside the repository under the large-file policy.
No manuscript snapshot directory is created.

\section{Outstanding analytical tasks}\label{sec:outstanding}

The proved chain ends at a uniform exterior enclosure and a collection
of localized and averaged moment estimates. Further progress requires
new quantitative input at identifiable points in that chain. The
following tasks are stated in the order of their immediate dependence.

\subsection{Acquire or test the every-center signed moment}

The nearest open positive task is to prove
\[
 Q(X)\le4/66^2\qquad\hbox{for every sufficiently large real }X,
\]
or a different complete boundary estimate below the Green margin in
\eqref{eq:margin-81}. Localization makes a uniformly strict cap on
$\|M_X\|_P$ sufficient after paying its norm error. The estimates in
Section~\ref{sec:means} allow rare large centers and therefore do not
complete this task. The missing quantity is a signed correlation with
the actual logarithmic phases, including nearby indices.

The complementary diagnostic task is the fixed short-window witness
$A_U(Y)\ge5/66^2$ along an unbounded sequence. Theorem~\ref{thm:quartic}
shows that one genuine alignment does not carry enough weighted mass.
Theorem~\ref{thm:sparse} further shows that any proposed mechanism
confined to $o(U^{1/64})$ of these logarithmic cells changes the full
moment by only $o(1)$. A new lower mechanism must control a more
distributed portion of the signed form, including its interaction
with the complement. It is not enough to display a constructive
subsum or to replace the complement by a growing upper allowance.
No such distributed mechanism is established in this manuscript.

\subsection{Extend the carrier and arithmetic control to deeper collars}

At the exact center on the curve,
\begin{equation}\label{eq:carrier-half}
 \Re s_{t_X}(X+\ii y)
 =\frac34+\frac y2+
          \frac{\log(1+(y/X)^2)}{8L_X}.
\end{equation}
The Euler-product margin is uniform for each fixed $y>1/2$.
At $y=1/2$, the remaining positive correction tends to zero;
below $1/2$ the real part is eventually below one. Thus the particular
elementary reciprocal-zeta proof used here loses its uniform margin
near $1/2$. This is a limitation of this carrier estimate, not a
theorem that every possible analytic carrier must fail.

An iteration needs a nonvanishing normalization or another signed
formulation on the actual deeper domains, together with complete
arithmetic error estimates there. The scale-invariant geometry of
Section~\ref{sec:geometry} does not supply either. Every stage must
also be available at the height under consideration: an unbounded
sequence of stage onsets would not justify infinitely many steps at
a fixed height.

\subsection{Pay the common-time compact region and the near-zero-time barrier}

The exterior theorem leaves $|x|<R(t)$, where $R(t)=4\pi\ee^{1/t}$.
At a fixed final time $T$, a direct zero exclusion on that full
compact complement would complete a global strip of the same width.
A different route is Polymath's three-part barrier criterion
\cite[Proposition 3.3]{Polymath}. To specify the missing work, fix
$0<y_0<1$, $0<T<1/2$ with $y_0^2+2T<1$, and set
$d_0=\sqrt{1-y_0^2}$. Its three zero-free regions are
\begin{align*}
 &t=0,\quad0\le x\le X_0,\quad
       \sqrt{y_0^2+2T}\le y\le1;\\
 &t=T,\quad x\ge X_0+d_0,\quad
       y_0\le y\le\sqrt{1-2T};\\
 &0\le t\le T,\quad X_0\le x\le X_0+d_0,\quad
       \sqrt{y_0^2+2(T-t)}\le y\le\sqrt{1-2t}.
\end{align*}
For $y_0=567/1000$, $T\le t_*$ and
$X_0+d_0\ge R(T)$, our theorem supplies the final exterior.
The initial region and the portion of the intermediate barrier with
$t<1/L_x$ remain outside its guarantee. In particular this includes
$t=0$. A complete certificate by either route would give a global
strip at $T$, and strip contraction would then imply
$\Lambda\le T+y_0^2/2$. No such certificate is supplied here.

The available onset is symbolic. Its sufficient construction already
requires $L_*\ge\kappa^{-1}=200000$, before the other constants are
evaluated. This describes the conservative proof recipe, not the true
first height at which the enclosure holds. A useful finite-height
bridge would require substantially sharper effective constants and
overlap with separately established compact control. The earlier
global-use audit's numerical comparison at width $4/7$ is specific
to that width~\cite{BB16}; it is not automatically a comparison for
$567/1000$.

\subsection{Justify any time-zero passage or iteration}

For every fixed finite $x$, the condition $|x|\ge4\pi\ee^{1/t}$
fails when $t$ is sufficiently close to zero. Hence the exterior
theorem cannot simply be sent to $t\downarrow0$ at that height.
Similarly, the conditional $9/16$ result is a curve statement;
it has no proved common-time extension. A limiting argument must
provide control on the relevant fixed compact sets, or a separate
uniform time-space propagation theorem with all boundary payments.

Finally, the complete corrected initial expansion and the joined
analytical chain merit specialist review independent of the shared
LLM source lineage. This would validate the foundations and the
new synthesis at a different evidence level. It would not by itself
fill the open moment, carrier or compact-region estimates. Those
remain the concrete mathematical tasks needed to turn the present
exterior progress into a global or iterative result.

\clearpage
\begingroup\small\raggedright
\begin{thebibliography}{99}
\bibitem{Polymath} D. H. J. Polymath.
\emph{Effective approximation of heat flow evolution of the Riemann
$\xi$ function, and a new upper bound for the de Bruijn--Newman constant}.
arXiv:1904.12438v2, 4 August 2019.
\url{https://arxiv.org/html/1904.12438v2}.
Theorem 1.5(i), Theorem 3.2 and Proposition 3.3 are the passages used.

\bibitem{MontgomeryVaughan} H. L. Montgomery and R. C. Vaughan.
\emph{Multiplicative Number Theory II: Primes and Sieves}, Chapter 16,
``Exponential Sums I: Van der Corput's Method.'' Cambridge, 2026.
Author-hosted text:
\url{https://personal.science.psu.edu/rcv4/571s25/montgomery-vaughanII.pdf}.
Definitions 16.1--16.2, (16.55)--(16.58), Theorems 16.23 and 16.25.
The investigation's source-use records 02 and 03 retain the checked
passage and version provenance.

\bibitem{DLMF} NIST Digital Library of Mathematical Functions.
Sections 20.7 (theta transformations), 25.2 (ordinary zeta series and
Euler product), 25.10(ii) (Riemann--Siegel expansion), and 5.11
(sectorial gamma asymptotics). \url{https://dlmf.nist.gov/}.

\bibitem{HDAComplex} E. Baker, with LLM assistance.
\emph{Fixed-cutoff complex remainder and buffered signed-time crossover}.
HDA Ansatz Note04, 25 October 2026. Repository path:
\path{papers/investigations-5/01_heat_density_asymptotics/ansatz/notes/04_COMPLEX_STRIP_REMAINDER_AND_BUFFERED_CROSSOVER_20261025.md}.
CS3, CS16--17 and their calibration inputs are inherited here.

\bibitem{HDAAssessment} E. Baker, with LLM assistance.
\emph{Complex strip and buffered crossover assessment}.
HDA Ansatz Review03, 25 October 2026. Repository path:
\path{papers/investigations-5/01_heat_density_asymptotics/ansatz/reviews/03_COMPLEX_STRIP_AND_BUFFERED_CROSSOVER_ASSESSMENT_20261025.md}.

\bibitem{BB06} E. Baker, with LLM assistance.
\emph{E01 r1---complete outer estimate and one inward step}.
BB Note06, 13 November 2026. In the present investigation's
\path{notes/06_GNP_BB_E01_COMPLETE_OUTER_ESTIMATE_20261113.md}.

\bibitem{BB12} E. Baker, with LLM assistance.
\emph{E05 r1---common-time width-$4/7$ tail}.
BB Note12, 13 November 2026.
\path{notes/12_GNP_BB_E05_COMMON_TIME_TAIL_20261113.md}.

\bibitem{BB16} E. Baker, with LLM assistance.
\emph{B04 global consumer audit for the common-time tail}.
BB Note16, 13 November 2026.
\path{notes/16_GNP_BB_B04_GLOBAL_CONSUMER_AUDIT_20261113.md}.

\bibitem{BB18} E. Baker, with LLM assistance.
\emph{E07 r1---five exponent pairs give width $567/1000$}.
BB Note18, 13 November 2026.
\path{notes/18_GNP_BB_E07_FIVE_PAIR_WIDTH_20261113.md}.

\bibitem{BB15} E. Baker, with LLM assistance.
\emph{B04: a signed Green acquisition on a fixed-depth collar}.
BB Note15, 13 November 2026.
\path{notes/15_GNP_BB_B04_SIGNED_GREEN_ACQUISITION_20261113.md}.

\bibitem{BB23} E. Baker, with LLM assistance.
\emph{GNP-BB-E09---localization and resonator means of the exact Poisson tail}.
BB Note23, 13 November 2026.
\path{notes/23_GNP_BB_E09_LOCALIZATION_AND_RESONATOR_MEANS_20261113.md}.

\bibitem{BB25} E. Baker, with LLM assistance.
\emph{B06: paid short-window freezing and one actual-phase packet test}.
BB Note25, 13 November 2026.
\path{notes/25_GNP_BB_B06_SHORT_WINDOW_AND_ACTUAL_PACKET_20261113.md}.

\bibitem{BBReview11} E. Baker, with LLM assistance.
\emph{E11 additional-agent review of the B06 short-window estimates}.
Review11, 13 November 2026.
\path{reviews/11_GNP_BB_E11_ADDITIONAL_AGENT_REVIEW_20261113.md}.
The explicitly dispatched reviewer used \texttt{gpt-6.1-sol}/\texttt{ultra}.

\bibitem{RepoResults} Shifted zeta positivity project.
\emph{Result summaries}, repository \path{result-summaries/README.md};
the relevant heat-flow cluster is
\path{result-summaries/11_newman_heat_flow_reductions_and_certificates.md}.
Read at the 13 November 2026 research state.

\bibitem{RepoAvoid} Shifted zeta positivity project.
\emph{Patterns to avoid}, repository \path{avoid/README.md}.
Read at the 13 November 2026 research state.

\bibitem{RepoAvoidPatterns} Shifted zeta positivity project.
Specific patterns: \path{avoid/01_restating_the_target.md},
\path{avoid/02_shape_only_arguments_and_the_controls.md},
\path{avoid/05_same_model_reviews_and_shared_premises.md},
\path{avoid/07_scoped_exclusions_read_as_universal.md}, and
\path{avoid/12_term_by_term_matching_and_separate_bounds.md}.

\bibitem{RepoControls} Shifted zeta positivity project.
\emph{Collision controls and the screening they impose}, repository
\path{result-summaries/12_collision_controls_and_screening.md}.
This catalog routes to the original constructions and their reviews;
the premise screen in this manuscript is scoped to those classes.
\end{thebibliography}

\endgroup
\end{document}
